\documentclass[11pt]{article}
\usepackage[T1]{fontenc}
\usepackage{lmodern,microtype}
\usepackage[margin=1in]{geometry}
\usepackage{amsmath,amssymb,amsthm,mathtools,bm}
\usepackage{booktabs,longtable,array,enumitem}
\usepackage{xcolor}
\usepackage[colorlinks=true,linkcolor=blue!45!black,citecolor=blue!45!black,urlcolor=blue!45!black]{hyperref}
\usepackage[nameinlink,noabbrev]{cleveref}
\hypersetup{pdftitle={Shape-Constrained Discovery-Rate Inversion},
 pdfauthor={Christopher D. Long},
 pdfsubject={Sampling calculus, recovery geometry, and finite-sample calibration},
 pdfkeywords={species sampling, moment problems, missing mass, optimal recovery}}
\usepackage{fancyhdr,aliascnt}
\pagestyle{fancy}\fancyhf{}
\fancyhead[L]{\small Shape-constrained discovery-rate inversion}
\fancyhead[R]{\small Revised draft v0.2}
\fancyfoot[C]{\thepage}
\setlength{\headheight}{14pt}
\setlength{\emergencystretch}{2em}
\setlist{nosep,leftmargin=*}
\numberwithin{equation}{section}
\newtheorem{theorem}{Theorem}[section]
\newaliascnt{proposition}{theorem}
\newtheorem{proposition}[proposition]{Proposition}
\aliascntresetthe{proposition}
\newaliascnt{lemma}{theorem}
\newtheorem{lemma}[lemma]{Lemma}
\aliascntresetthe{lemma}
\newaliascnt{corollary}{theorem}
\newtheorem{corollary}[corollary]{Corollary}
\aliascntresetthe{corollary}
\theoremstyle{definition}
\newaliascnt{definition}{theorem}
\newtheorem{definition}[definition]{Definition}
\aliascntresetthe{definition}
\newaliascnt{example}{theorem}
\newtheorem{example}[example]{Example}
\aliascntresetthe{example}
\theoremstyle{remark}
\newaliascnt{remark}{theorem}
\newtheorem{remark}[remark]{Remark}
\aliascntresetthe{remark}
\newcommand{\E}{\mathbb E}
\newcommand{\Pp}{\mathbb P}
\newcommand{\R}{\mathbb R}
\newcommand{\N}{\mathbb N}
\newcommand{\ind}{\mathbf1}
\newcommand{\Prob}{\mathcal P}
\newcommand{\Ccal}{\mathcal C}
\newcommand{\Scal}{\mathsf S}
\newcommand{\Tcal}{\mathsf T}
\newcommand{\Pois}{\operatorname{Pois}}
\newcommand{\Var}{\operatorname{Var}}
\newcommand{\Cov}{\operatorname{Cov}}
\newcommand{\supp}{\operatorname{supp}}
\newcommand{\osc}{\operatorname{osc}}
\newcommand{\TV}{d_{\mathrm{TV}}}
\newcommand{\Span}{\operatorname{span}}
\newcommand{\diag}{\operatorname{diag}}
\newcommand{\conv}{\operatorname{conv}}
\newcommand{\dd}{\,\mathrm d}
\newcommand{\ip}[2]{\langle #1,#2\rangle}
\DeclareMathOperator*{\argmax}{arg\,max}
\title{\textbf{Shape-Constrained Discovery-Rate Inversion}\\[5pt]
\large A sampling calculus, recovery geometry,\\and finite-sample statistical calibration}
\author{Christopher D. Long}
\date{November 12, 2026\\[3pt]\small Revised research draft, version 0.2}
\begin{document}
\maketitle
\begin{abstract}
Partial-population inference includes unseen-species prediction, missing-mass estimation, diversity estimation, capture--recapture, and ordered-support estimation. We organize these problems through sampling operations, target-directed recovery, and statistical calibration. A feature-incidence calculus treats coverage counts and maxima together: thinning generates a polynomial whose expansions at zero and one encode collisions and occupancy frequencies, while Poissonization converts discovery differences into derivatives. Convex moment geometry yields sharp relaxed bounds and contact certificates. Two canonical measures bound an entire discovery-rate curve from normalization, a rate, and its derivative; even derivative blocks give higher-order envelopes. We characterize the exceptional Hill orders by two one-dimensional families and distinguish relaxed measures from integer-realizable populations. Joint per-species cumulants give convex regions for a population measure and a realized or future random target at a single confidence level. Direct prediction of future discoveries remains exponentially calibratable even when the full Good--Toulmin error is not; variance, fourth-moment, and truncation alternatives are given. A perspective transformation makes projected Chernoff optimization a minimax equality on compact convex classes. Exact optimality is asserted only for specified models and criteria, including a support-clipped finite tank interval rule with a rational primal--dual certificate. Classical components are attributed, and the stated specializations are proved without claiming a universal estimator compiler or general statistical minimax theorem.
\end{abstract}
\noindent\textbf{Keywords.} Species sampling; thinning; moment problems; Chebyshev systems; missing mass; confidence regions; population realizability; optimal recovery.
\medskip
\noindent\textbf{Status.} Internally audited mathematical draft. Computational checks are supplementary, not proofs or formal verification. Novelty and independent peer review remain separate tasks.
\clearpage
\begingroup\footnotesize
\tableofcontents
\endgroup
\clearpage

\section{Introduction and scope}\label{sec:intro}
The organizing question is not which estimator to use, but which population functionals can be recovered from a specified experiment and what certifies that recovery. An experiment includes the retained information: labels, repetitions, order, batch membership, and capture histories can have different inferential content. Our construction has three layers:
\[
\text{sampling}\ \longrightarrow\ \text{recovery}\ \longrightarrow\ \text{calibration}\ \longrightarrow\ \text{optimality}.
\]
We keep four distinctions throughout: expected versus realized targets; exact moment information versus sampled information; positive mixing measures versus actual populations; and sharp projection of a confidence region versus globally optimal confidence procedures.

\subsection{Assumptions are target-specific}\label{subsec:assumptions}
An actual population is a finite or countable vector $p_i>0$ with $\sum_i p_i=1$. No positive minimum probability is imposed globally. The sample-reduction identities use exchangeability; per-species cumulant factorization uses independent Poisson counts. The compact envelope and integer-realizability results in \cref{sec:envelopes,sec:realizability} explicitly assume a known interval $0<a\le p_i\le b\le1$. This restriction ensures finiteness and compactness for support size and negative-order diversity. More general rare-species restrictions are possible, but are not interchangeable without proof.

In contrast, the finite-horizon discovery kernels, missing-mass kernels, and finite-coefficient joint cumulants extend continuously to zero. Their relaxed calibration can use all of $\Prob([0,1])$; an atom at zero in that relaxation is not an actual species. Unrestricted finite honest upper support bounds are impossible by \cref{prop:infiniteupper}. Minimum-mass support estimation, such as the class $p_i\ge1/k$ studied by Wu and Yang \cite{WuYang2019}, is therefore a distinct statistical problem rather than a harmless technical specialization.

\subsection{Historical position}
Good's frequency formulas and Good--Toulmin prediction are foundational \cite{Good1953,GoodToulmin1956}. Moment bounds for cataloging problems go back to Harris \cite{Harris1959}; Efron and Thisted used extrapolation and linear programming for unseen vocabulary \cite{EfronThisted1976}. Chao's richness inequality and higher-order moment approaches are prior art \cite{Chao1984,DaleySmith2016}. Discovery-rate representations of diversity profiles also precede this work \cite{ChaoJost2015}.

Poisson thinning and marking are standard \cite{LastPenrose2017}. Moment duality and extremal representations belong to classical convex and Chebyshev-system theory \cite{BertsimasPopescu2005,KarlinStudden1966,Pinelis2011}. Optimal recovery and statistical testing duality already connect approximating a target with constructing estimators and lower bounds \cite{Donoho1994,PolyanskiyWu2026,PolyanskiyWu2019}. Modern unseen-species prediction, occupancy concentration, and sequential calibration supply further baselines \cite{Orlitsky2016,BenHamou2017,Howard2021,Wasserman2020}. Eriksson's September 2026 version explicitly treats prediction, bounded-set incidence data, and optimality questions \cite{Eriksson2026}. The published general theory of Polyanskiy and Wu \cite{PolyanskiyWu2026} is distinguished from their combined preprint with unseen-population applications \cite{PolyanskiyWu2019}. The present draft should not be read as claiming the first unified statistical treatment of these topics.

Our aim is an explicit and reproducible synthesis across observation mechanisms. The most developed specializations are the common discovery-rate envelopes, their population-realizability restriction, and their interface with exact cumulant calibration. We do not establish a universal estimator compiler or a general minimax theorem for arbitrary partial-observation models.

\subsection{Main results}
The sampling results are \cref{thm:mobius,thm:differences,thm:covariance}. The recovery results are \cref{thm:momentdual,thm:minimax,thm:twonode,thm:evenblocks}. Joint calibration is based on \cref{thm:jointcalibration}, with explicit kernels in \cref{sec:errors}; population-only calibration is its specialization in \cref{thm:calibration}. The contact classification and Chernoff minimax identity are \cref{thm:intermediatehill,thm:chernoffminimax}. Exact model-specific optimality appears in \cref{thm:tankupper,prop:finitegame}.

The common Hill-diversity ordering holds for all real $q\le2$ and $q\ge3$ on a positive compact rate interval, while $2<q<3$ requires a sign test or the contact optimization in \cref{thm:intermediatehill}. Full unsmoothed Good--Toulmin extrapolation beyond the observed horizon can have no nonzero exponential moments (\cref{prop:gtmgf}). This limits exponential calibration of that estimator's error, not direct calibration of the future discovery count. The statistical sections retain both joint-region and explicit split-budget alternatives, without asserting uniform superiority of either.

\section{Sampling units, features, and information}\label{sec:model}
A sampling unit has a mark $z$ in a measurable space $\mathcal Z$. A feature is a measurable subset $A\subseteq\mathcal Z$; a unit reveals $A$ if its mark belongs to $A$. Let $w_A\ge0$ be a known weight. For a sample $S=(z_1,\ldots,z_n)$, set
\begin{equation}\label{eq:features}
 X_A(S)=\sum_{j=1}^n\ind_{\{z_j\in A\}},\qquad
 D(S)=\sum_Aw_A\ind_{\{X_A(S)>0\}},\qquad
 \Phi_r(S)=\sum_Aw_A\ind_{\{X_A(S)=r\}}\quad(r\ge1).
\end{equation}
Thus $D=\sum_{r\ge1}\Phi_r$. Feature memberships and weights must be computable from the recorded sample. Sample occurrences are distinguished even when their marks coincide.

Unless otherwise stated, feature calculations are for finitely many features. They extend to countable families when sums and expectations are justified by nonnegativity or absolute integrability. Extrapolation and covariance formulas require their own absolute-integrability conditions; they are not consequences of first-moment finiteness alone.

For species sampling the features are singleton labels. For a batch observation, a feature may mean that a particular species occurs anywhere in that batch. For a known phylogeny, a feature can be the descendant set of an edge, with edge length as weight. For positive integer serial numbers, take $A_m=\{z:z>m\}$ and $w_{A_m}=1$; then $D(S)=\max_jz_j$.

Under independent marks with law $Q$, write $q_A=Q(A)$. Under a finite population of $L$ units, write $a_A$ for its number of units revealing $A$. Avoidance probabilities are
\begin{equation}\label{eq:void}
 \Pp(X_A=0)=
 \begin{cases}
 (1-q_A)^n,&n\text{ independent units},\\
 e^{-tq_A},&\Pois(t)\text{ independent units},\\
 \binom{L-a_A}{n}/\binom Ln,&n\text{ uniform units without replacement}.
 \end{cases}
\end{equation}
We use $\binom uv=0$ when $v>u\ge0$. Features may overlap. Poissonization makes counts of disjoint mark categories independent, not arbitrary feature counts.

\begin{example}[The same mean is not the same experiment]\label{ex:dependence}
Two species with independent $\Pois(t/2)$ counts and two species whose counts are both one common $\Pois(t/2)$ variable have the same expected fingerprint and the same mean coverage $2(1-e^{-t/2})$. If $q=1-e^{-t/2}$, the discovery-count variances are respectively $2q(1-q)$ and $4q(1-q)$. In the second model the unit is a pair-revealing batch, with empty batches allowed to normalize the unit rate. First moments cannot replace the full sampling law.
\end{example}

\section{A calculus of sample reduction}\label{sec:sampling}
For a symmetric statistic $f$ defined on all subsamples, let
\begin{align}
 (\Scal_{n,m}f)(S)&=\binom nm^{-1}\sum_{I\subseteq[n]:\,|I|=m}f(S_I),\\
 (\Tcal_sf)(S)&=\sum_{I\subseteq[n]}s^{|I|}(1-s)^{n-|I|}f(S_I),\qquad 0\le s\le1.
\end{align}
These are conditional expectations under uniform subsampling and independent thinning. Direct composition gives
\begin{equation}\label{eq:semigroup}
 \Tcal_a\Tcal_b=\Tcal_{ab},\qquad
 \Scal_{n,m}\Scal_{m,k}=\Scal_{n,k},\qquad
 \Tcal_s=\sum_{m=0}^n\binom nm s^m(1-s)^{n-m}\Scal_{n,m}.
\end{equation}

\begin{theorem}[Subset-interaction coordinates]\label{thm:mobius}
For an occurrence-indexed subsample $B$, define
\[
 h_f(B)=\sum_{C\subseteq B}(-1)^{|B|-|C|}f(C).
\]
Then
\begin{align}
 f(S)&=\sum_{B\subseteq S}h_f(B),\\
 (\Tcal_sf)(S)&=\sum_{B\subseteq S}s^{|B|}h_f(B),\label{eq:mobiusthin}\\
 (\Scal_{n,m}f)(S)&=\sum_{B\subseteq S:\,|B|\le m}
 \frac{\binom m{|B|}}{\binom n{|B|}}h_f(B).\label{eq:mobiussub}
\end{align}
\end{theorem}
\begin{proof}
The first assertion is inclusion--exclusion, or Boolean-lattice M\"obius inversion. A fixed $k$-subset survives thinning with probability $s^k$ and uniform subsampling with probability $\binom mk/\binom nk$. Substitute the first identity and average.
\end{proof}
Thinning suppresses a $k$-unit interaction by $s^k$. Algebraic reversal consequently multiplies it by $s^{-k}$, explaining one source of inverse instability.

\begin{proposition}[Transfer under sample reduction]\label{prop:transfer}
For an exchangeable sequence whose first $n$ units exist,
\[
 \E_\theta(\Scal_{n,m}f)(S_n)=\E_\theta f(S_m).
\]
If $\ell(\cdot,T(\theta))$ is convex, then
\[
 \E_\theta\ell((\Scal_{n,m}f)(S_n),T(\theta))
 \le\E_\theta\ell(f(S_m),T(\theta)).
\]
\end{proposition}
\begin{proof}
A uniform $m$-subsample of an exchangeable $n$-sample has the $m$-sample law. Conditional Jensen proves the second assertion. The target is the same deterministic $T(\theta)$ on both sides.
\end{proof}
Known per-unit maps commute with these operations. An aggregate summary that discards unit boundaries need not do so. More generally, if an experiment is obtained from another by a known Markov kernel, any rule for the former can be simulated from the latter with exactly the same risk. This elementary transfer principle does not identify all sufficient summaries or all optimal rules.

\subsection{The coverage polynomial}
Define
\begin{equation}\label{eq:coveragepoly}
 P_S(s)=(\Tcal_sD)(S)=\sum_Aw_A[1-(1-s)^{X_A(S)}].
\end{equation}
Although the probabilistic definition uses $0\le s\le1$, the expression is a polynomial for every real $s$. Let
\[
 C_k(S)=\sum_Aw_A\binom{X_A(S)}k.
\]
Expanding at zero and one gives
\begin{align}
 P_S(s)&=\sum_{k\ge1}(-1)^{k+1}C_k(S)s^k,\label{eq:collisionexp}\\
 P_S(1+z)-P_S(1)&=\sum_{r\ge1}(-1)^{r+1}\Phi_r(S)z^r.\label{eq:fingerexp}
\end{align}
Consequently
\[
 P_S^{(k)}(0)=(-1)^{k-1}k!C_k(S),\qquad
 P_S^{(k)}(1)=(-1)^{k-1}k!\Phi_k(S).
\]
For $0\le s\le1$, $(-1)^{k-1}P_S^{(k)}(s)\ge0$. Collision statistics and occupancy frequencies are two expansions of the same sample polynomial.

Uniform subsampling instead yields the exact rarefaction identity
\begin{equation}\label{eq:rarefaction}
 (\Scal_{n,m}D)(S)=\sum_Aw_A
 \left[1-\frac{\binom{n-X_A(S)}m}{\binom nm}\right]
 =\sum_{k=1}^m(-1)^{k+1}\frac{\binom mk}{\binom nk}C_k(S).
\end{equation}
It is unbiased for expected $m$-sample coverage under any exchangeable design.

\begin{theorem}[Exchangeability identifies discovery differences]\label{thm:differences}
Let $F_n=\E D(S_n)$ and $\Delta F_n=F_{n+1}-F_n$. Whenever the exchangeable experiment exists through $n+r$,
\begin{equation}\label{eq:diffidentity}
 (-1)^{r-1}\Delta^rF_n
 =\frac{\E\Phi_r(S_{n+r})}{\binom{n+r}r},\qquad r\ge1.
\end{equation}
\end{theorem}
\begin{proof}
For one feature, inclusion--exclusion expresses the left side as the probability that it is absent in the first $n$ positions and present in every one of the next $r$ positions. Conditional on having exactly $r$ presences among $n+r$ positions, exchangeability makes each set of $r$ positions equally likely. Multiply by $w_A$ and sum.
\end{proof}
In particular $(-1)^{r-1}\Delta^rF_n\ge0$ on the available horizon and
\begin{equation}\label{eq:GTshift}
 \E\frac{\Phi_1(S_n)}n=F_n-F_{n-1}.
\end{equation}
For independent species observations this equals expected missing mass after $n-1$ observations, not after $n$. For batches it is an expected number or weight of newly revealed features, not necessarily a probability. Finite exchangeability does not imply an infinite completely monotone mixture representation.

\subsection{Poissonization and classical formulas}
For a sequence $a_n$, define $(\mathcal Pa)(t)=e^{-t}\sum_{n\ge0}t^na_n/n!$. Whenever termwise differentiation is justified,
\begin{equation}\label{eq:poissonintertwine}
 (\mathcal Pa)'=\mathcal P(\Delta a).
\end{equation}
With independent units and $S_t$ denoting a $\Pois(t)$ sample,
\begin{equation}\label{eq:poissonfeatures}
 F(t)=\E D(S_t)=\sum_Aw_A(1-e^{-tq_A}),\qquad
 (-1)^{r-1}F^{(r)}(t)=\frac{r!}{t^r}\E\Phi_r(S_t).
\end{equation}
Marginal Poisson counts suffice; disjoint features are not needed for these mean identities. The sample-size shift disappears because of \eqref{eq:poissonintertwine}. The standard Poisson deletion and marking identities give an equivalent derivation \cite{LastPenrose2017}.

For ordinary species, write $p=(p_i)$, $\sum_i p_i=1$, and
\begin{equation}\label{eq:mu}
 \mu_p=\sum_i p_i\delta_{p_i},\qquad
 r_p(t)=F'(t)=\int e^{-tx}\dd\mu_p(x).
\end{equation}
The function $r_p$ is completely monotone. Its log-convexity follows from
\[
 (\log r_p)''(t)=\Var_{\rho_t}(x)\ge0,
 \qquad \rho_t(\dd x)=e^{-tx}\mu_p(\dd x)/r_p(t).
\]

Differentiating \eqref{eq:coveragepoly} at $s=1$ gives the Poissonized Good--Turing formula $\Phi_1/t$, unbiased for $r_p(t)$. Evaluating at $s=1+u/t$ gives Good--Toulmin:
\begin{equation}\label{eq:GT}
 \widehat U_{t,u}=\sum_{r\ge1}(-1)^{r+1}(u/t)^r\Phi_r(S_t),\qquad
 \E\widehat U_{t,u}=F(t+u)-F(t).
\end{equation}
For finitely many features the expectation identity holds for every finite $u\ge0$. For ordinary countable species it also holds: if $K\sim\Pois(t)$ is total count, the absolute statistic is bounded by $K\max(1,u/t)^K$, which has finite expectation. The classical formulas and their regularization are discussed in \cite{Good1953,GoodToulmin1956,Orlitsky2016}.

The symmetry-based uniqueness theorem in \cite[Proposition 11, Definition 12, and Theorem 14]{Eriksson2026} is an important distinction from this algebraic derivation. In the near-future regime $u/t\le1$, the specified class combines the synchronous-expeditions property (additivity across disjoint populations) with atemporality (dependence on the effort ratio and data, not the absolute effort). Within that class Good--Toulmin is the unique rule with rate better than the null estimator under the paper's criterion. This is not uniqueness among all relabeling-invariant estimators. Effort-dependent regularization or restricted population models need separate comparisons. We cite the verified September 9, 2026 version; its prediction and incidence results are prior art, not consequences claimed for the first time here.

For fixed-size samples, the corresponding identity is instead
\[
 \E P_{S_n}(s)=\sum_{m=0}^n\binom nm s^m(1-s)^{n-m}F_m.
\]
Replacing $t$ by $n$ in \eqref{eq:GT} is not an exact fixed-size extrapolation argument.

Collision coefficients provide unbiased power sums:
\begin{equation}\label{eq:powers}
 \E\frac{C_k(S_n)}{\binom nk}=\sum_i p_i^k\quad(n\ge k),\qquad
 \E\frac{k!C_k(S_t)}{t^k}=\sum_i p_i^k.
\end{equation}
Nonlinear transformations, such as inverse Simpson diversity, do not preserve unbiasedness.

For $R(t)=\sum_i e^{-tp_i}$, $m_1=r_p(t)$ and $m_2=-r_p'(t)$, positivity gives
\[
 0\le\sum_i e^{-tp_i}(1-cp_i)^2=R(t)-2cm_1+c^2m_2.
\]
Optimizing $c$ proves
\begin{equation}\label{eq:Chao}
 R(t)\ge\frac{r_p(t)^2}{-r_p'(t)}
 =\frac{(\E\Phi_1)^2}{2\E\Phi_2}.
\end{equation}
This is the population moment inequality behind Chao-type richness bounds, not a confidence statement for a ratio of observed counts \cite{Chao1984,DaleySmith2016}.

\subsection{Multiple future surveys}
For independent future windows of efforts $u_1,\ldots,u_k$, let $J_k$ be expected feature weight absent initially and present in every future window. Then
\begin{equation}\label{eq:multifuture}
 J_k(t;\bm u)=\sum_Aw_Ae^{-tq_A}\prod_{j=1}^k(1-e^{-u_jq_A})
 =(-1)^{k+1}\Delta_{u_1}\cdots\Delta_{u_k}F(t).
\end{equation}
Apply the same finite differences to $P_{S_t}(1+v/t)$ to obtain an unbiased estimator. Additional absence windows shift the baseline effort in the target kernel. For example, absence initially and in a future window $w$, with presence in future windows $u,v$, gives coefficients
\[
 \frac{(-1)^r}{t^r}\{w^r-(w+u)^r-(w+v)^r+(w+u+v)^r\},\qquad r\ge2.
\]
These formulas generate mean estimators. Their variance and prediction coverage require separate analysis.

\subsection{Second-order information}
Write $q_{AB}=Q(A\cap B)$.
\begin{theorem}[Covariance calculus]\label{thm:covariance}
For a Poisson sample of independent units and finitely many features,
\begin{equation}\label{eq:covpoly}
 \Cov(P_{S_t}(s),P_{S_t}(v))
 =\sum_{A,B}w_Aw_B e^{-t(sq_A+vq_B)}(e^{tsv q_{AB}}-1)
\end{equation}
for all real $s,v$.
\end{theorem}
\begin{proof}
For a feature pair, split mark space into $A\setminus B$, $B\setminus A$, $A\cap B$, and the complement. Independent Poisson marking gives
\[
 \E(1-s)^{X_A}(1-v)^{X_B}=e^{-t(sq_A+vq_B-sv q_{AB})}.
\]
Subtract the product of marginal expectations and sum. The calculation remains valid for negative bases because all Poisson probability-generating series converge absolutely at finite arguments.
\end{proof}
There is also an unbiased covariance estimator. Let $N_{10},N_{01},N_{11}$ be the observed membership-cell counts for each ordered pair $A,B$. With $0^0=1$, set
\begin{align}\label{eq:covestimate}
 \widehat V_t(s,v)=\sum_{A,B}w_Aw_B(1-s)^{N_{10}}(1-v)^{N_{01}}
 \bigl( [(1-s)(1-v)]^{N_{11}}-(1-s-v)^{N_{11}}\bigr).
\end{align}
The cell probability-generating functions show that its expectation is \eqref{eq:covpoly}. A pair with $N_{11}=0$ contributes zero, so only observed co-occurring features are needed. It may be negative even at $s=v$, and unbiasedness does not provide a confidence bound.

\subsection{Two limitations of exact inversion}
\begin{proposition}[Fixed-size extrapolation obstruction]\label{prop:fixedobstruction}
For $m>n\ge1$, no statistic of $n$ independent observations is unbiased for expected $m$-sample species richness uniformly over the three-species probability simplex.
\end{proposition}
\begin{proof}
Any such expectation is a polynomial of degree at most $n$ in the probabilities. But $F_m(p)=3-\sum_{i=1}^3(1-p_i)^m$ has degree $m$ after $p_3=1-p_1-p_2$: the mixed leading coefficient of $p_1^{m-1}p_2$ is $-m$, which cannot cancel with the other two terms.
\end{proof}

\begin{proposition}[Long-horizon Good--Toulmin lacks exponential moments]\label{prop:gtmgf}
If $u/t>1$, the unsmoothed statistic \eqref{eq:GT} has finite moments of every finite positive order, but need not have any finite nonzero exponential moment. This failure already holds for a one-species population.
\end{proposition}
\begin{proof}
For the latter population, $X\sim\Pois(t)$ and the statistic is $0$ at $X=0$ and $(-1)^{X+1}a^X$ otherwise, where $a=u/t>1$. All power moments are bounded by $\E a^{kX}<\infty$. For $\lambda>0$, the odd-index terms in its exponential-moment series are
\[
 e^{-t}\frac{t^r}{r!}\exp(\lambda a^r),
\]
and do not tend to zero. For $\lambda<0$, use even indices. The finite-moment assertion for general species follows from the bound $K\max(1,a)^K$ above.
\end{proof}
This does not contradict \eqref{eq:covpoly}, which concerns second moments, not exponential moments of $P_{S_t}(s)$. It does restrict the prediction calibrations in \cref{sec:errors}.

\section{Recovery geometry and deterministic optimality}\label{sec:geometry}
Let $X$ be a compact metric space, $k_1,\ldots,k_d,g\in C(X)$, and $\nu\in\Prob(X)$. Write
\[
 A\nu=\left(\int k_1\dd\nu,\ldots,\int k_d\dd\nu\right),\qquad
 T_g(\nu)=\int g\dd\nu.
\]
The observable--target body is
\begin{equation}\label{eq:jointbody}
 \mathcal J_g=\{(A\nu,T_g(\nu)):\nu\in\Prob(X)\}
 =\conv\{(k_1(x),\ldots,k_d(x),g(x)):x\in X\}.
\end{equation}
It is compact and convex. For a feasible $m$, its vertical section is the closed interval
\[
 [L_g(m),H_g(m)]
 =\left[\min_{A\nu=m}T_g(\nu),\max_{A\nu=m}T_g(\nu)\right].
\]
Every intermediate value is attained in this relaxed class. $L_g$ is convex and $H_g$ concave. For multiple targets, their joint attainable body contains more information than a product of marginal intervals; its supporting directions reduce to scalar targets $\sum_jc_jg_j$.

\begin{theorem}[Moment duality]\label{thm:momentdual}
For every feasible $m$,
\begin{align}
 H_g(m)&=\inf_{c_0,c}\{c_0+c^\top m:c_0+c^\top k(x)\ge g(x)\ (x\in X)\},\label{eq:upperdual}\\
 L_g(m)&=\sup_{c_0,c}\{c_0+c^\top m:c_0+c^\top k(x)\le g(x)\ (x\in X)\}.
\end{align}
The measure extrema are attained. Dual optima are attained at relative-interior moment vectors; at boundary vectors only equality of the infimum or supremum values is asserted.
\end{theorem}
\begin{proof}
Integrate any pointwise certificate for weak duality. Conversely, the upper boundary of the compact convex body \eqref{eq:jointbody}, extended by $-\infty$ off its moment domain, is a closed concave function. It is the infimum of its affine majorants. An affine majorant on the moment body is equivalent to its inequalities at the generating points $k(x)$. This proves \eqref{eq:upperdual}; apply it to $-g$ for the lower bound. A finite concave function has a supporting affine majorant at every relative-interior point of its domain.
\end{proof}
These are standard moment-problem principles \cite{BertsimasPopescu2005}. The explicit distinction between value duality and certificate attainment prevents a boundary assumption from being silently inserted.

\subsection{Bounded errors and approximation}
Suppose only $\|A\nu-y\|\le\delta$ is known. For a nonempty feasible set,
\begin{align}
 H_{g,\delta}(y)
 &=\inf_c\left[c^\top y+\delta\|c\|_*+
       \sup_{x\in X}\{g(x)-c^\top k(x)\}\right],\label{eq:noisydual}\\
 L_{g,\delta}(y)
 &=\sup_c\left[c^\top y-\delta\|c\|_*+
       \inf_{x\in X}\{g(x)-c^\top k(x)\}\right].
\end{align}
Indeed, thicken $\mathcal J_g$ in its first coordinates by the radius-$\delta$ error ball and use \cref{thm:momentdual}; its support function is $\delta\|c\|_*$. Coordinatewise bounds $|e_j|\le\delta_j$ give the penalty $\sum_j\delta_j|c_j|$ instead.

\begin{theorem}[Modulus and optimal deterministic recovery]\label{thm:minimax}
Define
\[
 \omega_g(\epsilon)=\sup\{|T_g(\nu)-T_g(\eta)|:
                   \|A\nu-A\eta\|\le\epsilon\}.
\]
For deterministic errors bounded by $\delta$,
\begin{align}\label{eq:detminimax}
 \inf_{\widehat T}\sup_{\nu,\,\|e\|\le\delta}
 |\widehat T(A\nu+e)-T_g(\nu)|
 &=\tfrac12\omega_g(2\delta)\\
 &=\inf_{c_0,c}\left[
 \|g-c_0-c^\top k\|_\infty+\delta\|c\|_*\right].\nonumber
\end{align}
An affine rule attains the last value. In particular, universal exact recovery from $A\nu$ is possible if and only if $g\in\Span\{1,k_1,\ldots,k_d\}$.
\end{theorem}
\begin{proof}
Two moment vectors at distance at most $2\delta$ have a common possible observation, their midpoint. Any rule errs by at least half the target separation there. Conversely, the midpoint of the compatible target interval at each $y$ errs by at most half its width. This proves the first equality.

For the second, apply separating-hyperplane duality to the compact convex difference body
\[
 \{(A\nu-A\eta,T_g(\nu)-T_g(\eta)):\nu,\eta\in\Prob(X)\}.
\]
The dual value for half its maximal height over $\|z\|\le2\delta$ is
\[
 \inf_c\left[\delta\|c\|_*+\tfrac12\osc_{x\in X}(g(x)-c^\top k(x))\right].
\]
Choosing $c_0$ as the midpoint of the residual range converts half the oscillation into the uniform norm. This is the worst error of the affine rule $c_0+c^\top y$. A minimizer exists: for $\delta>0$ coefficient norms are bounded on sublevel sets; at $\delta=0$, use best approximation in a finite-dimensional closed subspace of $C(X)$. Exact recovery is equivalent to zero approximation distance.
\end{proof}
This is deterministic optimal recovery, not exact stochastic minimaxity for an arbitrary experiment. Its statistical relatives have substantial prior theory \cite{Donoho1994,PolyanskiyWu2019}.

\subsection{Contact sets and finite support}
If an attained upper certificate has residual
\[
 q(x)=c_0+c^\top k(x)-g(x)\ge0,
\]
then for a compatible $\nu$,
\[
 H_g(m)-T_g(\nu)=\int q\dd\nu.
\]
Thus every maximizing measure is supported on $\{q=0\}$. There is always an optimizer with at most $d+1$ atoms: an extreme moment-feasible measure cannot assign positive mass to $d+2$ disjoint measurable sets, because their $(d+1)$-dimensional vectors of mass and moments admit a signed, constraint-preserving perturbation. A continuous linear functional has an extreme optimizer on the compact feasible set.

On $[a,b]$, an extended Chebyshev space of dimension $d+2$ has at most $d+1$ zeros, counted with multiplicity, for each nonzero member. If $1,k_1,\ldots,k_d,g$ span such a space and the residual certificate is nonzero, its nonnegativity makes interior zeros even. Therefore
\begin{equation}\label{eq:contactcount}
 2n_{\rm interior}+n_{\rm endpoints}\le d+1.
\end{equation}
The nonzero-certificate qualification is essential; an exactly representable target has zero residual and imposes no such support restriction. General moment-extreme and Chebyshev principles are summarized in \cite{Pinelis2011}.

\section{Canonical simultaneous discovery envelopes}\label{sec:envelopes}
Assume throughout this section
\[
 0<a<b\le1,\qquad t>0,\qquad \nu\in\Prob([a,b]),\qquad
 r_\nu(s)=\int e^{-sx}\dd\nu(x).
\]
The class is a convex relaxation of actual species populations. Fix
\begin{equation}\label{eq:twomoments}
 r_\nu(t)=r_0,\qquad -r_\nu'(t)=r_1.
\end{equation}
Define the tilted probability measure and its mean by
\[
 \rho(\dd x)=\frac{e^{-tx}}{r_0}\nu(\dd x),\qquad \beta=r_1/r_0.
\]
Then \eqref{eq:twomoments}, including normalization, is equivalent to
\begin{equation}\label{eq:tiltmoments}
 \int\dd\rho=1,\qquad \int x\dd\rho=\beta,
 \qquad \int e^{tx}\dd\rho=1/r_0.
\end{equation}

\begin{lemma}[Feasibility and two endpoint representations]\label{lem:feasible}
The data are feasible exactly when $a\le\beta\le b$ and
\begin{equation}\label{eq:feasibility}
 e^{t\beta}\le1/r_0\le
 \frac{b-\beta}{b-a}e^{ta}+\frac{\beta-a}{b-a}e^{tb}.
\end{equation}
For strict inequalities and $a<\beta<b$, unique nodes $d\in(a,\beta)$ and $c\in(\beta,b)$ give representations
\begin{align}
 \rho_a&=\frac{c-\beta}{c-a}\delta_a+
          \frac{\beta-a}{c-a}\delta_c,\label{eq:rhoa}\\
 \rho_b&=\frac{b-\beta}{b-d}\delta_d+
          \frac{\beta-d}{b-d}\delta_b,\label{eq:rhob}
\end{align}
with the exponential moment in \eqref{eq:tiltmoments}. Put $\nu_a(\dd x)=r_0e^{tx}\rho_a(\dd x)$ and $\nu_b(\dd x)=r_0e^{tx}\rho_b(\dd x)$.
\end{lemma}
\begin{proof}
Jensen and the endpoint chord of $e^{tx}$ give necessity. Mixing the point mass at $\beta$ with the endpoint distribution of mean $\beta$ proves sufficiency. The exponential expectation for \eqref{eq:rhoa} is
\[
 e^{ta}+(\beta-a)\frac{e^{tc}-e^{ta}}{c-a},
\]
strictly increasing in $c$ by strict convexity. The corresponding expectation for \eqref{eq:rhob} is strictly decreasing in $d$. Their endpoint values are precisely the two extremes in \eqref{eq:feasibility}, proving existence and uniqueness. Untilting preserves all required moments.
\end{proof}
On the lower feasibility boundary the only measure is $\delta_\beta$; on the upper boundary the tilted endpoint measure is unique. Interpret both canonical measures there as that common measure.

\subsection{An elementary interpolation proof}
Let $\partial$ denote differentiation in $x$ and set
\begin{equation}\label{eq:operator}
 \mathcal L_t=\partial(\partial+t)^2
 =\partial^3+2t\partial^2+t^2\partial.
\end{equation}
Its nullspace is $\Span\{1,e^{-tx},xe^{-tx}\}$.

\begin{lemma}[Supporting interpolation]\label{lem:interpolation}
For $g\in C^3([a,b])$, if $\mathcal L_tg\le0$, then every measure compatible with \eqref{eq:twomoments} satisfies
\begin{equation}\label{eq:genorder}
 \int g\dd\nu_b\le\int g\dd\nu\le\int g\dd\nu_a.
\end{equation}
If $\mathcal L_tg\ge0$, both inequalities reverse. For strict sign and interior moment data the extremizers are unique.
\end{lemma}
\begin{proof}
Put $f=e^{tx}g$. Then $(e^{-tx}f'')'=\mathcal L_tg$. Interpolate $f$ by $Q=A+Bx+Ce^{tx}$ at $a,c$, with derivative matching at $c$. Uniqueness follows because a nonzero such $Q$ has at most two zeros with multiplicities: $Q''=Ct^2e^{tx}$ has constant sign when $C\ne0$, and the remaining case is affine.

Let $h=f-Q$. For strict negative operator sign, $(e^{-tx}h'')'<0$. A function with four zeros counted with multiplicity would, by two applications of Rolle followed by one weighted application, force a zero of $(e^{-tx}h'')'$. Thus $h$ has no additional zero beyond $a,c,c$. Also $h''(c)<0$: otherwise weighted monotonicity makes $h''>0$ to the left of $c$, and $h'(c)=0$ implies $h(a)>h(c)$, a contradiction. Hence $h<0$ off its interpolation nodes. This proves $Q\ge f$.

For interpolation at $d,d,b$, the same argument gives $h''(d)>0$ and $h>0$ off the nodes, producing a minorant. Integrate these inequalities under the tilted measures; their expectations depend only on \eqref{eq:tiltmoments}. Undoing the tilt gives \eqref{eq:genorder}. For a weak sign replace $f$ by $f+\epsilon x^2$, since $(e^{-tx}(f+\epsilon x^2)'')'$ is strictly negative, and pass to the limit. Apply the result to $-g$ for the opposite sign. Strictness forces equality measures to be supported on the two interpolation nodes; their moments then fix the weights.
\end{proof}

\begin{theorem}[Two-measure envelopes]\label{thm:twonode}
Let $\nu_-:=\nu_b$ and $\nu_+:=\nu_a$. For every compatible $\nu$,
\begin{equation}\label{eq:rateenvelope}
 r_{\nu_-}(s)\le r_\nu(s)\le r_{\nu_+}(s)\qquad(s\ge0).
\end{equation}
The same pair attains all the displayed bounds simultaneously. For interior data and $s>0$ with $s\ne t$, both extremizers are unique.
\end{theorem}
\begin{proof}
For $g_s(x)=e^{-sx}$,
\begin{equation}\label{eq:ratesign}
 \mathcal L_tg_s=-s(t-s)^2e^{-sx}\le0.
\end{equation}
Use \cref{lem:interpolation}. At $s=0,t$ the values are fixed by the constraints. Elsewhere the sign is strict.
\end{proof}
This is an explicit specialization of principal-representation theory. The argument above does not depend on numerical optimization.

Integrating \eqref{eq:rateenvelope} against any nonnegative measure gives sharp simultaneous bounds whenever the integral is finite. In particular it covers
\begin{equation}\label{eq:targets}
 F(v)=\int_0^v r(s)\dd s,\quad
 U(v,u)=\int_v^{v+u}r(s)\dd s,\quad
 R(v)=\int_v^\infty r(s)\dd s,\quad
 N=\int_0^\infty r(s)\dd s.
\end{equation}
These equal the corresponding abundance-measure functionals with kernels $(1-e^{-vx})/x$, $e^{-vx}(1-e^{-ux})/x$, $e^{-vx}/x$, and $1/x$. The assumption $a>0$ ensures finiteness, including for relaxed measures.

\subsection{Diversity, with the exceptional interval audited}
Define
\[
 D_q(\nu)=\left(\int x^{q-1}\dd\nu\right)^{1/(1-q)}\quad(q\ne1),
 \qquad D_1(\nu)=\exp\left(-\int\log x\dd\nu\right).
\]
These coincide with Hill diversities on actual populations. Negative orders are mathematically well-defined here because $a>0$; this extension does not assert robustness to arbitrarily rare species.

\begin{corollary}[Hill-diversity envelopes]\label{cor:hill}
For every real $q\in(-\infty,2]\cup[3,\infty)$,
\begin{equation}\label{eq:hille}
 D_q(\nu_-)\le D_q(\nu)\le D_q(\nu_+).
\end{equation}
For $2<q<3$, let $\tau_q=\sqrt{q-2}-(q-2)$. The order \eqref{eq:hille} holds if $ta\ge\tau_q$ and reverses if $tb\le\tau_q$. In particular, $ta\ge1/4$ guarantees \eqref{eq:hille} for every finite real $q$.
\end{corollary}
\begin{proof}
Direct differentiation gives
\begin{align}
 \mathcal L_t x^{q-1}
 &=(q-1)x^{q-4}\bigl[(tx+q-2)^2+2-q\bigr],\label{eq:hillsign}\\
 \mathcal L_t(-\log x)&=-\frac{(tx-1)^2+1}{x^3}<0.
\end{align}
For $q\le2$ or $q\ge3$ the bracket in \eqref{eq:hillsign} is nonnegative. Its sign is also positive for $2<q<3$ above $tx=\tau_q$, and negative below. Use \cref{lem:interpolation}; for $q>1$, the Hill transformation reverses the power-sum order. Finally, $\max_{0<z<1}(\sqrt z-z)=1/4$.
\end{proof}
Thus one must not extend the common ordering to all diversity orders without a condition. For $p=(0.5,0.3,0.2)$, $[a,b]=[0.05,0.6]$, $t=0.3$, and $q=2.25$, the canonical lower-rate, actual, and upper-rate measures give, respectively,
\[
 D_q\approx2.5977111,\quad2.5952556,\quad2.5873518.
\]
The reversed order follows analytically from $tb=0.18<\tau_{2.25}=0.25$, not only from the numerical values.

\subsection{Sharp bounds when the Hill sign changes}\label{subsec:intermediate}
The common pair need not optimize $D_q$ when $2<q<3$ and $a<\tau_q/t<b$. Nevertheless the sharp relaxed problem remains one-dimensional in each of two exhaustive support families.

\begin{theorem}[Contact classification for intermediate Hill orders]\label{thm:intermediatehill}
Assume interior feasible data in \eqref{eq:twomoments}, $2<q<3$, and $ta<\tau_q<tb$. Both extrema of $\int x^{q-1}\dd\nu$ over the compatible probability measures are attained, and every extremizing measure has support contained in either
\[
 \{x,y\},\quad a\le x<\beta<y\le b,
 \qquad\text{or}\qquad
 \{a,z,b\},\quad a<z<b.
\]
Degenerate zero weights and endpoint nodes are allowed. Equivalently, with $c,d$ from \cref{lem:feasible}, both sharp values can be found by optimizing over two compact one-parameter families of tilted measures:
\begin{align}
 \rho_x&=\frac{y(x)-\beta}{y(x)-x}\delta_x+
          \frac{\beta-x}{y(x)-x}\delta_{y(x)},&&x\in[a,d],\label{eq:hilltwofamily}\\
 \rho_z&=w_a(z)\delta_a+w_z(z)\delta_z+w_b(z)\delta_b,
          &&z\in[d,c].\label{eq:hillthreefamily}
\end{align}
Here $y(x)\in[c,b]$ is the unique solution of
\begin{equation}\label{eq:hilltwoconstr}
 \frac{(y-\beta)e^{tx}+(\beta-x)e^{ty}}{y-x}=1/r_0,
\end{equation}
and $(w_a,w_z,w_b)$ is the unique solution of
\begin{equation}\label{eq:hillthreeconstr}
 \begin{pmatrix}1&1&1\\a&z&b\\e^{ta}&e^{tz}&e^{tb}\end{pmatrix}
 \begin{pmatrix}w_a\\w_z\\w_b\end{pmatrix}
 =\begin{pmatrix}1\\\beta\\1/r_0\end{pmatrix}.
\end{equation}
For either tilted family the objective is
$r_0\int e^{tx}x^{q-1}\dd\rho$. Since $1/(1-q)<0$, its minimum gives the maximum Hill diversity and conversely.
\end{theorem}
\begin{proof}
By \cref{thm:momentdual}, an optimal certificate lies in $\Span\{1,e^{-tx},xe^{-tx}\}$; denote it by $Q$. For an upper bound, let $h=Q-g\ge0$, where $g(x)=x^{q-1}$. For a lower bound use $h=g-Q\ge0$. By \eqref{eq:hillsign}, $\mathcal L_th$ has exactly one zero in $(a,b)$, and this zero is simple. Three successive weighted applications of Rolle's theorem imply that $h$ has at most four zeros counted with multiplicity: five zeros would force at least two zeros of $\mathcal L_th$. Analyticity precludes an interval of zeros. An interior contact of a nonnegative residual has multiplicity at least two. Thus
\[
 2n_{\rm interior}+n_{\rm endpoints}\le4.
\]
Complementary slackness confines every extremizer to its contact set. Three distinct nodes can therefore occur only with both endpoints; otherwise there are at most two. A single node or an endpoint-only pair would put the moment vector on a feasibility boundary, excluded here.

For a two-node tilted measure the mean fixes its weights. For fixed $x<\beta$, strict convexity of $e^{t\cdot}$ makes the left side of \eqref{eq:hilltwoconstr} strictly increasing in $y$. Its limit at $y=\beta$ is $e^{t\beta}<1/r_0$, while its value at $y=b$ is at least $1/r_0$ exactly for $x\le d$. Its value at $y=c$ is at most $1/r_0$ for $x\ge a$. This proves the unique feasible branch on $[a,d]$, with $y(a)=c$ and $y(d)=b$.

For three nodes, the matrix in \eqref{eq:hillthreeconstr} is nonsingular by strict convexity. Cramer's rule, or comparison with the endpoint chords of $e^{tx}$, gives $w_z>0$, $w_a\ge0$ exactly when $z\ge d$, and $w_b\ge0$ exactly when $z\le c$. Hence its feasible range is $[d,c]$. Both families are continuous and compact, and their common endpoints reproduce the canonical representations. Optimizing them is exhaustive by the contact argument. Boundary moment data instead have the unique representation of \cref{lem:feasible}.
\end{proof}
A sampled grid or a local scalar optimizer is only a numerical approximation to these sharp extrema. Certified computation still requires global one-dimensional optimization or a pointwise dual certificate. The theorem concerns the positive-measure relaxation; integer multiplicities can alter its optimizer.

Retaining normalization also strengthens the two-moment Chao bound. In tilted form
\[
 R(t)=r_0\int x^{-1}\dd\rho\ge r_0/\beta=r_0^2/r_1.
\]
The exact relaxed lower value is instead $r_0\int x^{-1}\dd\rho_b$. It is strictly larger for interior data, because $\rho_b$ is then nondegenerate and $x^{-1}$ strictly convex.

\subsection{Several sampling times and derivative parity}
\begin{lemma}[Sign of a real-rooted Hermite remainder]\label{lem:hermitesign}
Let $\mathcal L=\prod_{j=1}^{n}(\partial-\lambda_j)$ have real characteristic roots, counted with multiplicity, and let $\mathcal U=\ker\mathcal L$. Let $Q\in\mathcal U$ interpolate $g\in C^n([a,b])$ at nodes $z_i$ with positive multiplicities $r_i$ summing to $n$. If $\mathcal Lg$ has a strict constant sign, then, off the nodes,
\begin{equation}\label{eq:hermitesign}
 \operatorname{sgn}(g(x)-Q(x))
 =\operatorname{sgn}(\mathcal Lg)\,
   \operatorname{sgn}\prod_i(x-z_i)^{r_i}.
\end{equation}
\end{lemma}
\begin{proof}
The identity $(\partial-\lambda)h=e^{\lambda x}(e^{-\lambda x}h)'$ gives the weighted Rolle zero count, with multiplicities. Consequently every nonzero element of $\mathcal U$ has at most $n-1$ zeros, and $g-Q$ cannot have $n+1$ zeros when $\mathcal Lg$ has strict sign. Interpolation is unique.

For distinct interpolation nodes, write the error as the usual evaluation-determinant quotient for a basis $u_1,\ldots,u_n$ of $\mathcal U$, appending the column $g$ and the evaluation row at $x$. The base determinant has constant nonzero sign on ordered nodes. The appended determinant also has constant nonzero sign there, by the preceding zero count. Coalescing ordered nodes fixes their orientations from their Wronskians. Since $\mathcal L$ is monic,
\[
 W(u_1,\ldots,u_n,g)=W(u_1,\ldots,u_n)\mathcal Lg;
\]
this follows by eliminating the last derivative row with the differential equation for the $u_j$. Moving the evaluation row to its ordered location supplies the sign of $\prod_i(x-z_i)$. Confluent limits give repeated nodes; the zero count rules out vanishing away from them. This proves \eqref{eq:hermitesign}. The argument is the generalized divided-difference orientation used in Chebyshev interpolation \cite{KarlinStudden1966,Pinelis2011}.
\end{proof}

We use the following classical principal-representation fact: for a $(2M+1)$-dimensional extended Chebyshev moment space containing the constant function, each interior moment vector has two endpoint principal representations, each with $M+1$ atoms, one including $a$ and one including $b$. They extremize all targets in the corresponding generalized-convexity class. Karlin--Studden \cite{KarlinStudden1966} give the classical moment-space theory. In the terminology-free formulation of Pinelis \cite[Proposition 2(I) and (IIa)]{Pinelis2011}, the extrema are unique, independent of the added target within its oriented Chebyshev class, and have precisely these endpoint support bounds. The statement applies to boundary vectors through their unique representing measure. The orientation needed here is proved in \cref{lem:hermitesign}, not inferred from the words ``upper'' or ``lower'' principal representation. The following specialization verifies the relevant class explicitly.

\begin{theorem}[Even-block global envelopes]\label{thm:evenblocks}
Fix distinct positive times $t_j$ and positive even integers $m_j$, with $\sum_jm_j=2M$. Suppose normalization and the moments
\[
 r^{(k)}(t_j),\qquad 0\le k<m_j,
\]
form an interior feasible moment vector over $\Prob([a,b])$. There are two principal measures, each with $M+1$ atoms, such that
\[
 r_{\nu_-}(s)\le r_\nu(s)\le r_{\nu_+}(s)\qquad(s\ge0)
\]
for every compatible $\nu$. The upper-rate measure includes $a$; the lower-rate measure includes $b$. Boundary data have the corresponding unique, degenerate representation.
\end{theorem}
\begin{proof}
The measurement space is
\[
 \mathcal U=\Span\{1,x^ke^{-t_jx}:0\le k<m_j\}
 =\ker\mathcal L,\qquad
 \mathcal L=\partial\prod_j(\partial+t_j)^{m_j}.
\]
Its dimension is $n=2M+1$ and its characteristic roots are real. Weighted Rolle, as in \cref{lem:hermitesign}, proves the extended Chebyshev property. The cited principal-representation theorem supplies an endpoint-$a$ measure and an endpoint-$b$ measure, each with $M$ interior nodes and one endpoint, depending only on the moment vector.

For any $g$ with $\mathcal Lg<0$, interpolate at the endpoint-$a$ nodes with multiplicity one at $a$ and two at each interior node. The product in \eqref{eq:hermitesign} is $(x-a)\prod_{i=1}^M(x-c_i)^2>0$ off the nodes. Thus $g-Q_a<0$ and $Q_a$ is a majorant. At the endpoint-$b$ nodes the product is $(x-b)\prod_{i=1}^M(x-d_i)^2<0$, so $g-Q_b>0$ and $Q_b$ is a minorant. Integrating these certificates identifies the endpoint-$a$ measure as the maximizer and the endpoint-$b$ measure as the minimizer.

For $g_s=e^{-sx}$,
\begin{equation}\label{eq:multisign}
 \mathcal Lg_s=-s\prod_j(t_j-s)^{m_j}e^{-sx}.
\end{equation}
All multiplicities are even, so this is strictly negative except at $s=0$ and the constrained times, where $g_s\in\mathcal U$. The same two measures therefore bound all rates simultaneously. Boundary data have the unique representing measure supplied by the classical moment theorem.
\end{proof}
When a block multiplicity is odd, \eqref{eq:multisign} changes sign across that time; the principal measures exchange their roles. Their support patterns depend on the total dimension. Thus the even-block hypothesis is structural, not a numerical convenience. In particular the first $2M$ expected Poisson occupancy counts at one time, together with normalization, give the above $M+1$-node envelopes.

\subsection{A numerical illustration}\label{subsec:example}
For $p=(1/2,3/10,1/5)$, $t=3$, and $[a,b]=[0.05,0.6]$, the exact-model moments are $r_0\approx0.3432983052$ and $r_1\approx0.1143262749$. Solving the two scalar equations in \cref{lem:feasible} gives the following values.
\begin{center}
\begin{tabular}{lrrr}
\toprule
Target & Relaxed lower & Actual value & Relaxed upper\\
\midrule
Future discoveries, $u=2$ &0.504478&0.505417&0.507366\\
Unseen richness at $t=3$ &1.107071&1.178511&2.030768\\
Total support &2.926594&3.000000&3.855201\\
Inverse Simpson diversity &2.603353&2.631579&2.670775\\
\bottomrule
\end{tabular}
\end{center}
The Chao value is $1.030854$. These are optimization values for known moments, not confidence limits based on one sample. The double-precision implementation is a numerical solver, not an interval-arithmetic certificate.

\section{Actual populations and integer realizability}\label{sec:realizability}
A probability mixing measure $\nu=\sum_\ell w_\ell\delta_{x_\ell}$ on positive rates is induced by a finite population if and only if
\begin{equation}\label{eq:integrality}
 w_\ell/x_\ell\in\N\quad\text{for every atom},\qquad \sum_\ell w_\ell=1.
\end{equation}
The integer is the number of species at that level. Equivalently, an integer counting measure $\eta=\sum_\ell n_\ell\delta_{x_\ell}$ satisfies $\nu(\dd x)=x\eta(\dd x)$ and $\int x\dd\eta=1$. Thus an actual-population optimization is
\[
 \text{optimize }\int xg(x)\dd\eta(x),\quad
 \int xk_j(x)\dd\eta=m_j,\quad \int x\dd\eta=1,
 \quad\eta\text{ integer-valued}.
\]
There is no general reason for a relaxed principal representation to satisfy \eqref{eq:integrality}.

\begin{proposition}[A strict relaxation gap]\label{prop:relaxgap}
For rates in $[a,b]$ with $a>0$, the class of actual population measures is compact. If an exact-moment relaxed optimizer is unique and unrealizable, then its value is strictly separated from the actual-population optimum, provided the latter feasible set is nonempty.
\end{proposition}
\begin{proof}
There are at most $\lfloor1/a\rfloor$ species. For each possible number, the closed subset of $[a,b]^N$ on which coordinates sum to one is compact; its image under $p\mapsto\sum_i p_i\delta_{p_i}$ is compact in the weak topology. Take the finite union. Moment constraints are closed and the target continuous. An actual optimum exists, and equality with the relaxed value would identify its measure with the unique unrealizable optimizer.
\end{proof}

\begin{example}[Integrality collapses an interval]\label{ex:integergap}
Assume $p_i\in[1/4,1/2]$ and $\sum_i p_i^2=3/8$. In the relaxation, $\int x\dd\nu=3/8$ and $N=\int x^{-1}\dd\nu$. Jensen gives $N\ge8/3$, attained at $\delta_{3/8}$. The chord inequality $x^{-1}\le6-8x$ gives $N\le3$. All intermediate values are attainable by convexity. Actual $N$ is integral, hence equals $3$, and $(1/2,1/4,1/4)$ realizes it. The relaxed target interval is $[8/3,3]$, while the actual one is $\{3\}$.
\end{example}
The few-node results of \cref{sec:geometry,sec:envelopes} are not few-level theorems for every actual-population objective. This manuscript does not claim a universal two-level principle.

\section{Joint population--target calibration}\label{sec:calibration}
Return to an actual species population $p$ and independent counts $X_i\sim\Pois(tp_i)$. Fix an integer $d\ge1$, put $\Phi=(\Phi_1,\ldots,\Phi_d)$, and define
\[
 b_r(x)=e^{-tx}(tx)^r/r!,\qquad k_r(x)=b_r(x)/x,
 \qquad m_r=\E\Phi_r=\int k_r\dd\mu_p.
\]
The fingerprint coordinates are not independent. Independence across species, which is specific to this Poissonized model, is what permits the following factorization.

\begin{theorem}[An exact linear cumulant transform]\label{thm:cumulant}
For $\lambda\in\R^d$, define
\begin{equation}\label{eq:kappa}
 \kappa_\lambda(x)=\frac1x\log\left(1+\sum_{r=1}^d b_r(x)(e^{\lambda_r}-1)\right).
\end{equation}
Then
\begin{equation}\label{eq:cgf}
 \log\E_p e^{\lambda^\top\Phi}=\int\kappa_\lambda\dd\mu_p.
\end{equation}
The kernel has a continuous extension at zero, with value $t(e^{\lambda_1}-1)$.
\end{theorem}
\begin{proof}
One species contributes at most one coordinate unit to $\Phi$. Its exponential moment is the quantity inside the logarithm. Multiply across species and use $\sum_i f(p_i)=\int f(x)/x\dd\mu_p(x)$. The argument of the logarithm is strictly positive because it is a categorical exponential moment. Its expansion at zero gives the stated limit. For countably many species, $\kappa_\lambda$ is bounded on $[0,1]$, and $e^{\lambda^\top\Phi}$ is dominated by $e^{C K}$ for total Poisson count $K$, which justifies the limiting identity.
\end{proof}
Thus the log moment-generating function, not only the expected fingerprint, is a linear functional of the abundance measure. An arbitrary relaxed measure need not define a population law; \eqref{eq:cgf} is a distributional statement only at actual population measures.

\subsection{Joint transforms for additive random targets}\label{subsec:joint}
Let $Z=\sum_i Z_i\in\R^k$ be a realized or future random target, where the pairs $(X_i,Z_i)$ are independent across species and their per-species law depends on the population only through $x=p_i$. Define
\begin{equation}\label{eq:jointkernelgeneral}
 \kappa_{\lambda,\eta}(x)
 =\frac1x\log\E_x\exp\left(
    \sum_{r=1}^d\lambda_r\ind_{\{X=r\}}+\eta^\top Z_x\right).
\end{equation}
For finitely many species with finite exponential moments, factorization gives
\begin{equation}\label{eq:jointcgf}
 \log\E_p e^{\lambda^\top\Phi+\eta^\top Z}
 =\int\kappa_{\lambda,\eta}\dd\mu_p.
\end{equation}
For countably many species we require convergence of $Z$, convergence of the log-product, and uniform integrability of the exponentials of partial sums. These conditions, not additivity alone, justify \eqref{eq:jointcgf}. They are verified explicitly for the applications in \cref{sec:errors}. The vector formulation allows finitely many target coordinates to be calibrated together.

\begin{theorem}[Joint exponential calibration and projection]\label{thm:jointcalibration}
Let $K\subseteq[0,1]$ be compact, and assume the kernels \eqref{eq:jointkernelgeneral} for finitely many fixed directions $(\lambda_j,\eta_j)$ extend to continuous finite functions on $K$. Suppose \eqref{eq:jointcgf} holds for every actual population in the model. Choose fixed $w_j>0$ with $\sum_jw_j=1$. For a fixed, nonrandom, closed convex candidate class $\Ccal\subseteq\Prob(K)$ containing the true $\mu_p$, define
\begin{equation}\label{eq:jointregion}
 \mathcal J_\alpha(\Phi)=
 \left\{(\nu,z)\in\Ccal\times\R^k:
 \sum_jw_j\exp\left(\lambda_j^\top\Phi+\eta_j^\top z
             -\int\kappa_{\lambda_j,\eta_j}\dd\nu\right)
 \le\alpha^{-1}\right\}.
\end{equation}
This set is convex and
\begin{equation}\label{eq:jointcoverage}
 \Pp_p\{(\mu_p,Z)\in\mathcal J_\alpha(\Phi)\}\ge1-\alpha.
\end{equation}
Its projection onto $z$ is a convex prediction or confidence set for $Z$. For a scalar target its closed interval hull has coverage at least $1-\alpha$. Deterministic support restrictions, such as $0\le M_t\le1$, may be intersected before projection. For integer targets, intersect the resulting scalar interval with the integers. The level is a single $\alpha$; no population/error split is required.
\end{theorem}
\begin{proof}
At $(\mu_p,Z)$ each unweighted exponential in \eqref{eq:jointregion} has expectation one by \eqref{eq:jointcgf}; their weighted sum does too. Markov's inequality proves \eqref{eq:jointcoverage}. For the observed $\Phi$, every exponent is affine in $(\nu,z)$, giving joint convexity. Projection preserves convexity and inclusion of the true pair implies inclusion of the true target. If the reported projection is empty, a vacuous full-support set is a valid always-reported convention. Measurability holds in the compact continuous models here, with endpoints interpreted in the extended reals.
\end{proof}
A previously estimated population confidence region cannot be substituted for the fixed class $\Ccal$ at no coverage cost: its failure event must also be accounted for, or its calibration incorporated in the same joint construction. The construction is an application of e-value inversion and projection, not a new general validity principle \cite{Wasserman2020,Howard2021}. It does not assert conditional coverage given every fingerprint, bounded interval width, or optimal length. Directions, weights, future horizons, and target coefficients are predeclared; independent training permits conditional application, whereas selection on the evaluation sample requires a separate calibration argument. Actual-population multiplicities form a nonconvex subclass: restricting to it retains coverage, but its projection need not be an interval; the interval hull remains valid. There is no claim that joint calibration uniformly dominates every split-budget construction.

\subsection{Population-only regions and simultaneous projection}
A nonnegative statistic $E_\theta(X)$ satisfying $\E_\theta E_\theta\le1$ supplies a test by Markov's inequality. This is the usual e-value construction; likelihood-based versions and sequential developments are well established \cite{Wasserman2020,Howard2021}.

\begin{theorem}[Convex exponential calibration]\label{thm:calibration}
Choose finitely many directions $\lambda^{(j)}$ and weights $w_j>0$, $\sum_jw_j=1$, independently of the evaluation sample. For $\nu\in\Prob(K)$, with $K=[0,1]$ or a specified compact rate interval, set
\[
 E_\nu(\Phi)=\sum_jw_j\exp\left((\lambda^{(j)})^\top\Phi
                               -\int\kappa_{\lambda^{(j)}}\dd\nu\right).
\]
The convex set
\begin{equation}\label{eq:confset}
 \Ccal_\alpha(\Phi)=\{\nu:E_\nu(\Phi)\le1/\alpha\}
\end{equation}
contains the true $\mu_p$ with probability at least $1-\alpha$ whenever its rates belong to $K$.
An alternative valid convex region is the intersection of linear constraints
\begin{equation}\label{eq:linconf}
 \int\kappa_{\lambda^{(j)}}\dd\nu\ge
 (\lambda^{(j)})^\top\Phi-\log(1/(\alpha w_j)).
\end{equation}
\end{theorem}
\begin{proof}
This is \cref{thm:jointcalibration} with no random target coordinate, using \cref{thm:cumulant}. Each constraint in \eqref{eq:linconf} fails with probability at most $\alpha w_j$; use a union bound.
\end{proof}
The linear-constraint region is a relaxation of the mixture region, since a sum bounded by $1/\alpha$ bounds each weighted summand. It may be easier to optimize, but need not be equally efficient.

For any nonempty confidence region define
\begin{equation}\label{eq:projection}
 L_T(X)=\inf_{\nu\in\Ccal_\alpha(X)}T(\nu),\qquad
 H_T(X)=\sup_{\nu\in\Ccal_\alpha(X)}T(\nu).
\end{equation}
On the event that the true measure belongs to the region, all deterministic population-function bounds hold simultaneously. Selecting a target or deterministic horizon after observing the data is therefore valid when it only selects a projection of this same region. Selecting calibration directions on the evaluation sample is a different operation and is not justified by this argument. Empty feasible regions may be reported as rejection of the model/calibration constraints; for an always-reported interval, use a vacuous full-model bound on such outcomes.

\subsection{A mean-only calibration and the two-dimensional reduction}
Using $\log(1+z)\le z$ in \eqref{eq:kappa},
\begin{equation}\label{eq:poissondom}
 \log\E e^{\lambda^\top\Phi}\le\sum_{r=1}^d m_r(e^{\lambda_r}-1).
\end{equation}
This is exponential-moment domination by independent Poisson coordinates, not an assertion of independence or of Poisson fingerprint marginals.

Define the Poisson deviance
\[
 h(z,m)=z\log(z/m)-z+m,
\]
with $h(0,m)=m$, $h(0,0)=0$, and $h(z,0)=\infty$ for $z>0$. Optimizing the one-coordinate Chernoff inequality separately in its upper and lower tails gives
\begin{equation}\label{eq:deviance}
 \Pp\left(h(\Phi_r,m_r)\le\log(2d/\alpha)
                         \text{ for all }r\le d\right)\ge1-\alpha.
\end{equation}
Indeed, for a fixed true $m$, $h(z,m)$ is decreasing up to $m$ and increasing afterwards. Each tail beyond a level $c$ has probability at most $e^{-c}$; sum the $2d$ bounds. This monotonic tail inversion is why optimizing the Chernoff parameter here is valid without paying an uncountable union penalty.

For $d=2$, $r_0=m_1/t$ and $r_1=2m_2/t^2$. Explicitly, on a known interval $[a,b]$ with $a>0$, set
\begin{equation}\label{eq:gammaexplicit}
 \Gamma_\alpha(\Phi)=\left\{(r_0,r_1):
 \begin{array}{l}
 h(\Phi_1,tr_0)\le\log(4/\alpha),\\
 h(\Phi_2,t^2r_1/2)\le\log(4/\alpha),\\
 (r_0,r_1)\text{ satisfies \eqref{eq:feasibility}}
 \end{array}\right\}.
\end{equation}
This is precisely the $d=2$ deviance region, intersected with deterministic moment feasibility. Then for every target in the common envelope class,
\begin{equation}\label{eq:twodimconf}
 L_T=\inf_{\xi\in\Gamma_\alpha}T(\nu_-(\xi)),\qquad
 H_T=\sup_{\xi\in\Gamma_\alpha}T(\nu_+(\xi))
\end{equation}
is a simultaneous confidence construction. The inner infinite-dimensional optimization has been solved; the outer search is two-dimensional. Different targets can optimize at different $\xi$. Integrating pointwise extrema gives valid but not necessarily sharp projected integral bounds. Additional exact cumulant constraints generally destroy this particular two-moment reduction; they require a new recovery calculation.

\section{Realized targets, prediction, and stopping}\label{sec:errors}
A confidence interval for an expected quantity is not automatically an interval for the realized missing mass or a future discovery count. We first apply the joint theorem directly to those random targets, then give explicit and computationally convenient alternatives. All future efforts and predictor coefficients below are fixed independently of the evaluation sample.

\subsection{Direct calibration of missing mass}\label{subsec:jointmissing}
Let
\[
 M_t=\sum_i p_i\ind_{\{X_i=0\}},\qquad \widehat M_t=\Phi_1/t.
\]
For $Z_i=p_i\ind_{\{X_i=0\}}$, the joint kernel in \eqref{eq:jointkernelgeneral} is
\begin{equation}\label{eq:jointmissing}
 \kappa^M_{\lambda,\eta}(x)
 =\frac1x\log\left[
 1+\sum_{r=1}^d b_r(x)(e^{\lambda_r}-1)
       +e^{-tx}(e^{\eta x}-1)\right].
\end{equation}
The missing-mass and positive-count contributions are mutually exclusive for each species. Thus
\[
 \log\E_p e^{\lambda^\top\Phi+\eta M_t}
 =\int\kappa^M_{\lambda,\eta}\dd\mu_p.
\]
The kernel extends continuously to zero with value
\begin{equation}\label{eq:missinglimit}
 \kappa^M_{\lambda,\eta}(0)=t(e^{\lambda_1}-1)+\eta.
\end{equation}
For countably many species, $M_t\le1$ and $|\lambda^\top\Phi|\le\|\lambda\|_\infty K_t$, where $K_t\sim\Pois(t)$ is total initial count. These bounds justify uniform integrability and convergence in \eqref{eq:jointcgf}.

\begin{corollary}[Single-level realized-mass interval]\label{cor:jointmissing}
Use \eqref{eq:jointregion} with the kernel \eqref{eq:jointmissing}, candidate measures in $\Prob([0,1])$, and target coordinate $z\in[0,1]$. The projection onto $z$ is an interval with unconditional coverage at least $1-\alpha$ for $M_t$. No positive lower probability bound or nontrivial upper probability bound is required.
\end{corollary}
\begin{proof}
The preceding calculation verifies all hypotheses of \cref{thm:jointcalibration}; projection supplies the result. The interval depends only on the observed fingerprint and the fixed calibration choices, not on the unobserved species probabilities.
\end{proof}
This is not a conditional-coverage claim given every fingerprint, nor a minimum-length theorem. A poorly chosen finite set of calibration directions may give a vacuous interval. The region treats $z$ as a free candidate target coordinate; imposing additional genuine structural restrictions can narrow it, but inventing a deterministic relation between $z$ and $\nu$ cannot be justified by the expectation formula alone.

\subsection{A closed-form asymmetric alternative}\label{subsec:bernstein}
The error is a sum of independent centered terms
\[
 Z_i=t^{-1}\ind_{\{X_i=1\}}-p_i\ind_{\{X_i=0\}}.
\]
Since the two indicators are mutually exclusive,
\begin{equation}\label{eq:missingvariance}
 V:=\Var(\widehat M_t-M_t)
 =\sum_i(p_i/t+p_i^2)e^{-tp_i}
 =\frac{m_1+2m_2}{t^2}.
\end{equation}
Thus $(\Phi_1+2\Phi_2)/t^2$ is unbiased for this exact Poissonized error variance; that fact alone is not a confidence theorem.

There is a distinct fixed-size precedent. In Esty's normal-limit theorem, as recorded and extended by Zhang and Zhang \cite{Esty1983,ZhangZhang2009}, the studentizing expression for the fixed-$n$ missing-mass error is
\begin{equation}\label{eq:estyvariance}
 \frac{\Phi_1(1-\Phi_1/n)+2\Phi_2}{n^2}
 =\frac{\Phi_1+2\Phi_2}{n^2}-\frac{\Phi_1^2}{n^3}.
\end{equation}
It is part of an asymptotic fixed-sample theorem under its stated hypotheses, not the exact Poisson identity \eqref{eq:missingvariance}. Zhang and Zhang also give the corresponding Poisson variance scale \cite[Section 2]{ZhangZhang2009}. Nonasymptotic occupancy and missing-mass bounds are developed in \cite{BenHamou2017}.

\begin{theorem}[An asymmetric calibrated missing-mass interval]\label{thm:missingci}
Suppose $p_i\le b\le1$ and $V_U(X)\ge0$ exceeds $V$ with probability at least $1-\alpha_0$, for example by projection of a calibrated moment region. With $\ell=\log(2/\beta)$, the interval
\begin{equation}\label{eq:missingci}
 \left[
 \widehat M_t-\sqrt{2V_U\ell}-\frac{\ell}{3t},\quad
 \widehat M_t+\sqrt{2V_U\ell}+\frac{b\ell}{3}
 \right]\cap[0,1]
\end{equation}
covers $M_t$ with probability at least $1-\alpha_0-\beta$. The choice $b=1$ is always valid and imposes no additional population restriction.
\end{theorem}
\begin{proof}
We have the separate bounds $Z_i\le1/t$ and $-Z_i\le b$. For a centered $Z\le c$ and $\lambda\ge0$, the one-sided Bennett bound follows from
\[
 e^{\lambda z}\le1+\lambda z+
 \frac{z^2}{c^2}(e^{\lambda c}-1-\lambda c)\qquad(z\le c).
\]
Summing logarithms and using $e^u-1-u\le u^2/[2(1-u/3)]$ for $0\le u<3$ gives the sub-gamma bound with variance $V$ and scale $c/3$. Apply it to $\sum_iZ_i$ with $c=1/t$ and to $-\sum_iZ_i$ with $c=b$. Each tail has probability at most $e^{-\ell}$ at its respective radius $\sqrt{2V\ell}+c\ell/3$; this is the standard Bernstein inversion \cite[Chapter 2]{Boucheron2013}. Intersect the two tail events with $V_U\ge V$. No independence between the variance bound and the error is required. Countable sums follow by the same second-moment and exponential bounds on finite partial sums.
\end{proof}
This formula is an inexpensive explicit alternative to joint projection, not a claimed consequence with identical endpoints or a uniformly inferior procedure.

\subsection{Direct prediction of future discoveries}\label{subsec:jointfuture}
Let independent future increments have duration $u\ge0$, and set
\[
 Y_{t,u}=\sum_i\ind_{\{X_i=0\}}\ind_{\{X_i^{\rm future}>0\}},\qquad
 q_{t,u}(x)=e^{-tx}(1-e^{-ux}).
\]
The joint transform of $(\Phi,Y_{t,u})$ is
\begin{equation}\label{eq:jointfuture}
 \kappa^Y_{\lambda,\eta}(x)
 =\frac1x\log\left[
 1+\sum_{r=1}^d b_r(x)(e^{\lambda_r}-1)
          +q_{t,u}(x)(e^\eta-1)\right],
\end{equation}
with continuous extension
\[
 \kappa^Y_{\lambda,\eta}(0)
 =t(e^{\lambda_1}-1)+u(e^\eta-1).
\]
Indeed, new discoveries require initial count zero and are therefore mutually exclusive with all positive-count fingerprint events. For countably many species, $Y_{t,u}\le K_u^{\rm future}$, an independent $\Pois(u)$ total, supplies domination along with $K_t$.

\begin{corollary}[Single-level future-count prediction]\label{cor:jointfuture}
Use \cref{thm:jointcalibration} with \eqref{eq:jointfuture}, $\nu\in\Prob([0,1])$, and $z\ge0$. Project onto $z$, take its closed interval hull, and intersect with $\mathbb Z_{\ge0}$. The resulting set covers $Y_{t,u}$ with probability at least $1-\alpha$ for every finite predeclared $t>0,u\ge0$.
\end{corollary}
\begin{proof}
The finite continuous kernels and the Poisson-total domination just established verify \eqref{eq:jointcgf}. Apply the joint projection theorem.
\end{proof}
The interval need not be centered at Good--Toulmin and may be wide. Crucially, the target $Y_{t,u}$ has exponential moments even when the full unsmoothed Good--Toulmin error does not. \Cref{prop:gtmgf} therefore obstructs a particular error-based calibration, not direct exponential calibration of the future count. This distinction complements the prediction results in \cite{Orlitsky2016,Eriksson2026}.

\subsection{Finite-coefficient error calibration}\label{subsec:finiteerror}
Fix coefficients $c_1,\ldots,c_d$ and define
\[
 \widehat Y=\sum_{r=1}^dc_r\Phi_r,\qquad E_{t,u}=\widehat Y-Y_{t,u}.
\]
The joint fingerprint--error kernel is
\begin{equation}\label{eq:jointerror}
 \kappa^E_{\lambda,\eta}(x)
 =\frac1x\log\left[
 1+\sum_{r=1}^db_r(x)(e^{\lambda_r+\eta c_r}-1)
          +q_{t,u}(x)(e^{-\eta}-1)\right].
\end{equation}
Its limit at zero is $t(e^{\lambda_1+\eta c_1}-1)+u(e^{-\eta}-1)$. Since $|\widehat Y|\le\max_r|c_r|K_t$, all fixed-direction transforms are finite. Joint projection onto the error coordinate gives limits $e_L,e_U$ and the single-level prediction set
\begin{equation}\label{eq:jointpredictioninterval}
 [\widehat Y-e_U,\widehat Y-e_L]\cap\mathbb Z_{\ge0}
\end{equation}
with coverage at least $1-\alpha$. The deterministic support condition $\widehat Y-e\ge0$, expressed using the observed predictor, can also be imposed before projection.

Setting $\lambda=0$ in \eqref{eq:jointerror} gives the error-only kernel
\begin{equation}\label{eq:predictionkernel}
 \kappa^{\rm pred}_\eta(x)=\frac1x\log\left[
 1+\sum_{r=1}^db_r(x)(e^{\eta c_r}-1)
        +q_{t,u}(x)(e^{-\eta}-1)\right],
\end{equation}
and
\begin{equation}\label{eq:predictioncgf}
 \log\E_p e^{\eta(\widehat Y-Y_{t,u})}
 =\int\kappa^{\rm pred}_\eta\dd\mu_p.
\end{equation}
These formulas retain the dependence between predictor and target. They cannot be extended to the full long-horizon series solely on the basis of the existence of its power moments.

\begin{theorem}[Projected Chernoff prediction]\label{thm:prediction}
Let $\Ccal_{\alpha_0}(X)$ be a population confidence region. Define
\[
 b_\pm(\nu)=\inf_{\lambda>0}
 \frac{\log(1/\beta)+\int\kappa^{\rm pred}_{\pm\lambda}\dd\nu}{\lambda},
 \qquad
 B_\pm(X)=\max\left(0,\sup_{\nu\in\Ccal_{\alpha_0}(X)}b_\pm(\nu)\right).
\]
Then
\begin{equation}\label{eq:predictioninterval}
 [\widehat Y-B_+(X),\widehat Y+B_-(X)]\cap\mathbb Z_{\ge0}
\end{equation}
has joint frequentist prediction coverage at least $1-\alpha_0-2\beta$.
\end{theorem}
\begin{proof}
At the fixed true population, Chernoff bounds and a limiting minimizing sequence give $\Pp(\widehat Y-Y>b_+(\mu_p))\le\beta$ and its negative-tail analogue. To justify a nonattained infimum, apply the bound first to any threshold strictly above the infimum, then decrease that threshold to it. On the population-coverage event the projected radii are no smaller. A union bound proves the result. The true-population optimizer is not selected using the evaluation data.
\end{proof}
This alternative retains an error split but has a useful exact optimization property. It is not conditional coverage given every initial sample, and a data-selected family of future random targets requires additional joint calibration.

\begin{theorem}[Exact Chernoff minimax interchange]\label{thm:chernoffminimax}
Let $\Ccal\subseteq\Prob(K)$ be nonempty, compact and convex, with $K\subseteq[0,1]$ compact. For the finite-coefficient kernel above and $\ell>0$,
\begin{equation}\label{eq:chernoffminimax}
 \sup_{\nu\in\Ccal}\inf_{\lambda>0}
 \frac{\ell+\int\kappa^{\rm pred}_{\pm\lambda}\dd\nu}{\lambda}
 =\inf_{\lambda>0}
 \frac{\ell+\sup_{\nu\in\Ccal}\int\kappa^{\rm pred}_{\pm\lambda}\dd\nu}{\lambda}.
\end{equation}
The outer problem is convex after $v=1/\lambda$; its infimum need not occur at a finite positive $\lambda$. The result applies pointwise to every nonempty compact convex realization of a calibrated region.
\end{theorem}
\begin{proof}
Write $K_\nu(\lambda)=\int\kappa^{\rm pred}_{\pm\lambda}\dd\nu$. Every per-species log moment-generating function is convex, so $K_\nu$ is convex and finite. With $v=1/\lambda$ the objective becomes
\[
 F(\nu,v)=\ell v+vK_\nu(1/v).
\]
It is affine and continuous in $\nu$ for each $v>0$. It is convex and continuous in $v$, since it is a restriction of the perspective of $K_\nu$; differentiating gives
\[
 \frac{\mathrm d^2}{\mathrm dv^2}[vK_\nu(1/v)]
 =\frac{K_\nu''(1/v)}{v^3}\ge0.
\]
Sion's minimax theorem, with the measure variable compact and the positive half-line convex, gives $\sup_\nu\inf_vF=\inf_v\sup_\nu F$ \cite[p.~174]{Sion1958}. This is \eqref{eq:chernoffminimax}. Each inner population optimization is linear, and the supremum of the convex functions of $v$ is convex.
\end{proof}
For an integer-realizable population class, convexity can fail. Then only the general inequality $\sup\inf\le\inf\sup$ is assured, and the right side is a conservative bound. The nonfinite error exponential transforms in \cref{prop:gtmgf} also fall outside this theorem; its finite-order moment cumulants still exist. The equality is sharp for the projected Chernoff criterion, not necessarily for exact tail quantiles or globally optimal prediction intervals.

\subsection{Power-moment repairs beyond the observed horizon}\label{subsec:gtmoments}
The absence of exponential moments does not preclude valid prediction. Let $S=\widehat U_{t,u}-Y_{t,u}$ and $a=u/t$. Define
\begin{equation}\label{eq:gtevenmoments}
 v_{2j}(x)=e^{-tx}\bigl(e^{t a^{2j}x}-e^{-ux}\bigr),\qquad j\ge1.
\end{equation}

\begin{proposition}[Full Good--Toulmin prediction variance]\label{prop:gtpredictionvariance}
For every finite $t>0,u\ge0$,
\begin{equation}\label{eq:gtpredvar}
 \E S=0,\qquad
 V_{t,u}:=\Var S
 =\int\frac{v_2(x)}x\dd\mu_p(x)
 =\int\frac{e^{-tx}(e^{u^2x/t}-e^{-ux})}{x}\dd\mu_p(x).
\end{equation}
If $V_U\ge V_{t,u}$ with probability at least $1-\alpha_0$, then
\[
 [\widehat U_{t,u}-\sqrt{V_U/\beta},\widehat U_{t,u}+\sqrt{V_U/\beta}]
 \cap\mathbb Z_{\ge0}
\]
has prediction coverage at least $1-\alpha_0-\beta$.
\end{proposition}
\begin{proof}
For a species of probability $x$, the predictor is $(-1)^{r+1}a^r$ when its initial count is $r\ge1$ and zero otherwise. Its mean equals $q_{t,u}(x)$, the per-species future discovery mean. The two contributions have zero product. Hence the centered per-species error has second moment $v_2(x)$; more generally its $2j$th moment is $v_{2j}(x)$. Sum the independent centered errors. The variance kernel has finite limit $u^2/t+u$ at zero. Chebyshev and the event $V_U\ge V_{t,u}$ prove the interval claim.
\end{proof}
A valid $V_U$ is obtained by projecting a population region onto this kernel. It is not asserted to belong to the common positive-Laplace envelope class.

\begin{proposition}[Fourth-moment prediction bound]\label{prop:fourthmoment}
Let
\begin{equation}\label{eq:gtfourth}
 K_4=\int\frac{v_4(x)-3v_2(x)^2}{x}\dd\mu_p(x),\qquad
 M_4=K_4+3V_{t,u}^2.
\end{equation}
Then $\E S^4=M_4$. If $M_{4,U}\ge M_4$ with probability at least $1-\alpha_0$, the radius $(M_{4,U}/\beta)^{1/4}$, centered at $\widehat U_{t,u}$ and intersected with the nonnegative integers, has prediction coverage at least $1-\alpha_0-\beta$.
\end{proposition}
\begin{proof}
For finitely many independent centered species errors, fourth cumulants add and $\E S^4=\sum_i(\E S_i^4-3(\E S_i^2)^2)+3(\sum_i\E S_i^2)^2$. Use \eqref{eq:gtevenmoments}. For a countable population, finite partial errors converge in $L^4$: their absolute values are dominated by $K_t\max(1,a)^{K_t}+K_u^{\rm future}$, which has all finite power moments. The kernels have finite limits at zero, giving the stated integrals. Apply Markov to $S^4$ and intersect with the upper-bound event. Cumulants here are defined by moments (or characteristic-function derivatives), not by an analytic moment-generating function near zero.
\end{proof}
One can project a calibrated region onto the expression in \eqref{eq:gtfourth}, clamping the supremum below at zero. This is generally a nonlinear optimization: the squared variance term must not be described as a linear moment objective.

Fourth moments do not uniformly improve Chebyshev. Even for known $V,M_4>0$, the fourth-moment radius is smaller exactly when
\begin{equation}\label{eq:fourthcomparison}
 \beta\le V^2/M_4.
\end{equation}
For one species, $t=1$, $u=2$ and $\beta=0.05$, the respective radii are approximately $20.0179$ and $89.9211$; the crossover is approximately $0.0001228$. Higher moments generate further valid alternatives, but selecting their smallest data-dependent radius requires simultaneous error allocation or another selection argument. Taking a minimum of separately level-$\beta$ bounds without adjustment is not automatically level-$\beta$.

\subsection{A calibrated truncation alternative}\label{subsec:truncation}
For a predeclared integer $R\ge1$, let
\[
 \widehat U_R=\sum_{r=1}^R(-1)^{r+1}(u/t)^r\Phi_r.
\]
It equals the full statistic unless some species has count exceeding $R$. Write $F_R(z)=\Pp\{\Pois(z)\le R\}$.
\begin{proposition}[Exact truncation failure transform]\label{prop:truncation}
For every finite or countable actual population,
\begin{align}
 \delta_R(p)&:=\Pp_p\{\max_iX_i>R\}
 =1-\exp\left(\int\frac{\log F_R(tx)}x\dd\mu_p(x)\right),\label{eq:truncationexact}\\
 \delta_R(p)&\le
 \int\frac{1-F_R(tx)}x\dd\mu_p(x).\label{eq:truncationunion}
\end{align}
Both kernels extend continuously to zero (with value zero for $R\ge1$). If a prediction procedure based on $\widehat U_R$ has failure probability at most $\gamma$, recentering its interval at the full $\widehat U_{t,u}$ has failure probability at most $\gamma+\delta_R(p)$, provided the widths are unchanged.
\end{proposition}
\begin{proof}
Independence gives $1-\delta_R=\prod_iF_R(tp_i)$; take logarithms to obtain \eqref{eq:truncationexact}. The union bound gives \eqref{eq:truncationunion}. The Poisson tail is $O(x^{R+1})$ at zero, proving the extensions. Infinite products are justified by summability of the tail probabilities (at most $t/(R+1)$ by Markov); also the no-arrival event has probability $e^{-t}$, so the product is positive. The last claim holds because the original and recentered intervals agree on the no-truncation event.
\end{proof}
For a uniform guarantee, bound $\delta_R(p)$ uniformly over the declared model by a predetermined $\epsilon$. A calibrated population region can also support an explicit gate: report the recentered interval only when its projected truncation probability is at most $\epsilon$, and otherwise report the full nonnegative range. If that region fails with probability at most $\alpha_0$, the gated rule has failure probability at most $\gamma+\alpha_0+\epsilon$. Indeed, when the true $\delta_R>\epsilon$ the gate can open only on population-region failure; when $\delta_R\le\epsilon$, use \cref{prop:truncation}. This explains precisely how a data-dependent tail certificate is used. A cutoff $R$ selected on the evaluation data needs an additional simultaneous or independent-selection justification.

\subsection{Independent batches}
For $n$ independent units revealing total feature weight at most $B$ per unit, replacing one unit changes any $\Phi_r$ by at most $2B$. Bounded differences and a union bound give
\begin{equation}\label{eq:batchbound}
 \Pp\left( |\Phi_r-\E\Phi_r|\le
 B\sqrt{2n\log(2d/\alpha)}\text{ for all }r\le d\right)\ge1-\alpha.
\end{equation}
This baseline is the bounded-differences argument of \cite[Chapter 6]{Boucheron2013}; it requires no within-unit independence. It is not claimed to be sharp. The independent-species cumulant formula \eqref{eq:cgf} does not apply to arbitrary batches or to fixed-size multinomial sampling. Stronger incidence prediction and concentration results exist \cite{Eriksson2026}; their assumptions must be checked before importing them.

\subsection{Optional stopping}
Fixed-time coverage across all prediction horizons is not time-uniform coverage across observation times. One can allocate errors across predetermined observation times. Alternatively, under a sequential model, choose before observation $X_j$ a proper predictive density $q_j(\cdot\mid X_{<j})$ and form
\[
 E_{\theta,n}=\prod_{j=1}^n
 \frac{q_j(X_j\mid X_{<j})}{p_\theta(X_j\mid X_{<j})}.
\]
On the true-model support this is a nonnegative supermartingale; with matching absolute continuity it is a martingale. Ville's inequality gives
\[
 \Pp_\theta\{\sup_nE_{\theta,n}>1/\alpha\}\le\alpha.
\]
Inverting the tests and projecting yields confidence sequences \cite{Howard2021,Wasserman2020}. The fixed-time fingerprint calibrators in \cref{thm:calibration} are not automatically supermartingales in $t$.

\section{Risk optimization and statistical obstructions}\label{sec:risk}
\subsection{A noise-aware affine estimator program}
Fix coefficients before observing the evaluation sample, and estimate $T_g(\mu_p)$ by $\widehat T=a+c^\top\Phi$. Let
\[
 u_c(x)=\sum_rc_rb_r(x),\qquad
 v_c(x)=\frac{\sum_rc_r^2b_r(x)-u_c(x)^2}{x}.
\]
Independence across species gives the exact mean-squared error
\begin{equation}\label{eq:risk}
 \E_p(\widehat T-T_g)^2
 =\left[\int(a+c^\top k-g)\dd\mu_p\right]^2+\int v_c\dd\mu_p.
\end{equation}
The variance matrix $\diag(b(x))-b(x)b(x)^\top$ is positive semidefinite, being a categorical covariance matrix. Jensen therefore gives the uniform risk certificate
\begin{equation}\label{eq:risksup}
 \E_p(\widehat T-T_g)^2\le
 \sup_{x\in[a,b]}\left[(a+c^\top k(x)-g(x))^2+v_c(x)\right].
\end{equation}
The supremum of the algebraic expression \eqref{eq:risk} over all probability mixing measures equals the right side, attained at a point mass. Minimizing it over $a,c$ is a convex optimization problem. It is an exact optimization of a relaxed affine risk certificate, not a proof of minimax optimality over all actual populations and all estimators. Extra moment restrictions can be retained in a sharper program.

Data-dependent coefficient selection needs independent training or a uniform selection argument. This restriction does not apply to selecting a deterministic target projection from an already valid population confidence region, because its coverage event is simultaneous. General estimator/testing dualities provide a broader context \cite{PolyanskiyWu2019}.

\subsection{Lower bounds use observation laws}
For two populations $\theta_0,\theta_1$ with target separation $\Delta$, every estimator satisfies
\begin{equation}\label{eq:twopoint}
 \max_{j=0,1}\E_{\theta_j}|\widehat T-T(\theta_j)|
 \ge\frac\Delta2\left[1-\TV(P_{\theta_0},P_{\theta_1})\right].
\end{equation}
Indeed, integrate the pointwise triangle inequality against the common part of the two laws. If an interval $I$ has coverage $1-\alpha$ at both populations, then
\begin{equation}\label{eq:intervalLB}
 \E_{\theta_0}|I|\ge
 \Delta[1-2\alpha-\TV(P_{\theta_0},P_{\theta_1})]_+.
\end{equation}
Under $\theta_0$, the interval contains both target values with at least the bracketed probability; on that event its length is at least $\Delta$.

A moment-matching pair does not by itself control total variation of the full experiment. Least-favorable priors supported on actual populations can yield stronger lower bounds than a single pair \cite{PolyanskiyWu2019}. If $P_\theta$ is passed through a known Markov kernel, total variation decreases. Consequently the target modulus over difficult-to-distinguish populations can only increase under coarsening.

\begin{proposition}[Unrestricted upper support bounds must be infinite]\label{prop:infiniteupper}
Suppose $U(X)$ is an honest $(1-\alpha)$ upper confidence bound for support size over all finite discrete populations, using a fixed finite independent sample or fixed finite Poisson effort. At a one-species population,
\[
 \Pp\{U=\infty\}\ge1-\alpha.
\]
\end{proposition}
\begin{proof}
Compare $p=(1)$ with $q_K=(1-\epsilon_K,\epsilon_K/K,\ldots,\epsilon_K/K)$ having $K$ additional species, with $\epsilon_K\downarrow0$. A coupling gives observation-law total variation at most $n\epsilon_K$ for $n$ observations and $t\epsilon_K$ at Poisson effort $t$. Honesty at $q_K$ implies under $p$ that $\Pp\{U\ge K+1\}\ge1-\alpha-o(1)$. The events decrease to $\{U=\infty\}$.
\end{proof}
A known positive minimum probability changes this problem fundamentally. Complete monotonicity alone does not exclude arbitrarily rare species.

\section{Ordered-support inference and exact decision games}\label{sec:tank}
Take a uniform sample of $k$ distinct integers from $\{1,\ldots,N\}$ and order it as $Y_{(1)}<\cdots<Y_{(k)}$, with $Y_{(0)}=0$. Under upper-tail features $A_m=\{z:z>m\}$,
\begin{equation}\label{eq:gaps}
 D=Y_{(k)},\qquad \Phi_r=Y_{(k-r+1)}-Y_{(k-r)}\quad(1\le r\le k).
\end{equation}
Thus feature singletons become the top spacing. The $k+1$ empty gaps are exchangeable, so estimating the final empty gap by the average observed empty gap gives $(k+1)Y_{(k)}/k-1$. Standard distributional and sufficiency results for this model are discussed in \cite{ONeill2025}.

Let $M=Y_{(k)}$. Its law is
\begin{equation}\label{eq:tanklaw}
 \Pp_N(M=m)=\frac{\binom{m-1}{k-1}}{\binom Nk},\qquad
 \Pp_N(M\le m)=\frac{\binom mk}{\binom Nk}\quad(k\le m\le N).
\end{equation}

\begin{proposition}[Triangular inversion]\label{prop:tankunbiased}
For every finite-valued function $\psi$ on $\{k,k+1,\ldots\}$,
\begin{equation}\label{eq:tankinverse}
 h_\psi(m)=\psi(m)+\frac{m-k}{k}[\psi(m)-\psi(m-1)],\quad m>k,
 \qquad h_\psi(k)=\psi(k)
\end{equation}
is the unique minimum-variance unbiased estimator of $\psi(N)$.
\end{proposition}
\begin{proof}
Unbiasedness is equivalent to
\[
 \sum_{m=k}^N h_\psi(m)\binom{m-1}{k-1}=\psi(N)\binom Nk.
\]
Subtract the equation at $N-1$ to obtain \eqref{eq:tankinverse}. The same triangular subtraction shows completeness of $M$. Conditional on $M=m$, the other $k-1$ sampled labels form a uniform subset of $\{1,\ldots,m-1\}$, independent of $N$, establishing sufficiency. Conditional averaging then proves the minimum-variance assertion. Each fixed $N$ has finite observation support, so all the expectations used are finite.
\end{proof}
Labels are essential. If only the ordinary multiplicity fingerprint of this without-replacement sample is retained, it is always $(k,0,\ldots)$ and supplies no information about $N$ beyond $N\ge k$.

\begin{theorem}[Pointwise smallest monotone upper limit]\label{thm:tankupper}
For $0<\alpha<1$, define
\begin{equation}\label{eq:tankCI}
 U_\alpha(m)=\max\left\{N\ge m:\frac{\binom mk}{\binom Nk}>\alpha\right\}.
\end{equation}
Then $\Pp_N(N\le U_\alpha(M))\ge1-\alpha$. Among nondecreasing, nonrandomized integer upper confidence bounds with this coverage, $U_\alpha$ is pointwise smallest.
\end{theorem}
\begin{proof}
The retained set is nonempty and finite. To verify monotonicity, extend $F_N(m)=\Pp_N(M\le m)$ by $F_N(m)=1$ for $m\ge N$. The sets $\{N\ge k:F_N(m)>\alpha\}$ are nested as $m$ increases; their maxima are precisely $U_\alpha(m)$. Hence $U_\alpha$ belongs to the declared nondecreasing class. Noncoverage occurs exactly when $F_N(M)\le\alpha$, and the discrete probability integral transform gives $\Pp_N(F_N(M)\le\alpha)\le\alpha$. If a competing nondecreasing bound $V$ has $V(m)<U_\alpha(m)$, set $N=U_\alpha(m)$. Every $M\le m$ then gives $V(M)<N$, while $\Pp_N(M\le m)>\alpha$, contradiction.
\end{proof}
The strict inequality in \eqref{eq:tankCI} is needed for the pointwise-minimal statement. Replacing it by $\ge$ remains conservative but can be larger at equality cases. This is one-sided optimality in a declared class, not optimality of every two-sided tail inversion.

For tag--retag with independent uniform captures of sizes $a,b$, overlap has law
\[
 \Pp_N(J=j)=\frac{\binom aj\binom{N-a}{b-j}}{\binom Nb}.
\]
Exact tail inversion is consequently available. The construction alone does not prove a shortest two-sided interval or derive every conventional overlap estimator as a unique optimum.

\subsection{A finite decision-theoretic certificate}
For finite parameter, observation, and interval-action sets, let $d_{m,I}$ be the probability of reporting interval $I$ at observation $m$. To minimize worst expected length, solve
\begin{align}\label{eq:game}
 \text{minimize }&R,\nonumber\\
 \sum_I d_{m,I}&=1,\quad d_{m,I}\ge0,\nonumber\\
 \sum_mP_N(m)\sum_{I\ni N}d_{m,I}&\ge1-\alpha\quad\text{for every }N,\\
 \sum_mP_N(m)\sum_I|I|d_{m,I}&\le R\quad\text{for every }N.\nonumber
\end{align}
This is a linear program optimizing coverage and length jointly. A dual certificate consists of risk weights $\lambda_N\ge0$, $\sum_N\lambda_N=1$, coverage weights $\gamma_N\ge0$, and numbers $v_m$ satisfying
\begin{equation}\label{eq:gamedual}
 v_m\le\sum_NP_N(m)[\lambda_N|I|-\gamma_N\ind_{\{N\in I\}}]
 \quad\text{for every }m,I.
\end{equation}
Every feasible rule then has worst expected length at least
\begin{equation}\label{eq:dualvalue}
 (1-\alpha)\sum_N\gamma_N+\sum_mv_m.
\end{equation}
This follows by multiplying \eqref{eq:gamedual} by $d_{m,I}$, summing, and applying coverage. Finite linear-programming duality supplies equality at an optimum.

\begin{proposition}[An exactly solved restricted tank game]\label{prop:finitegame}
For $k=3$, $N\in\{3,\ldots,12\}$, confidence level $19/20$, and actions consisting of all integer-endpoint intervals in $[3,12]$ plus the empty interval, allowing randomization, the minimax expected length is
\begin{equation}\label{eq:697}
 R_*=\frac{697}{140}.
\end{equation}
Length is upper endpoint minus lower endpoint; empty and singleton intervals have zero length.
\end{proposition}
\begin{proof}
An explicit feasible rule is given in \cref{tab:rule}. Its intervals are clipped to $[M,12]$, the feasible parameter range after observing $M$. For any interval $I$, membership of the true $N$ in $I$ is identical to membership in $I\cap[M,12]$ almost surely, while length cannot increase. This clipping improves the earlier untrimmed minimax rule at non-worst-case parameters without changing its optimal maximum risk. It does not by itself establish admissibility. Direct rational calculation gives coverage $19/20$ at every parameter and the risks in \cref{tab:risks}, whose maximum is $697/140$.

For an analytic lower certificate, choose
\[
 \lambda_N=\ind_{\{N=7\}},\qquad
 \gamma_N=\frac{\binom N3}{35},\qquad
 w_m=\frac{\binom{m-1}2}{35},\qquad
 v_m=\begin{cases}-w_m,&m\le7,\\-(13-m)w_m,&m>7.\end{cases}
\]
For $m\le N$, $\gamma_NP_N(m)=w_m$. If $m\le7$, the right side of \eqref{eq:gamedual} is
\[
 w_m\bigl[|I|-\#(I\cap\{m,\ldots,12\})\bigr]\ge-w_m.
\]
For $m>7$ it is $-w_m\#(I\cap\{m,\ldots,12\})\ge-(13-m)w_m$. The empty action also satisfies the inequalities. Finally,
\[
 \sum_{N=3}^{12}\gamma_N=143/7,\qquad
 \sum_{m=3}^{12}v_m=-101/7,
\]
so \eqref{eq:dualvalue} is $(19/20)(143/7)-101/7=697/140$. This matches the exhibited rule.
\end{proof}
Since $M$ is sufficient, allowing other information in the unordered sampled subset does not improve this randomized finite decision problem: condition any such rule on $M$. The result makes no assertion about unbounded $N$, nonrandomized rules, or a different loss.

\section{What the combined theory establishes}\label{sec:discussion}
For a specified sampling model, the calculus yields observable kernels and, where available, exact error transforms. Geometry characterizes their compatible population functionals. Calibration attaches a coverage event to constraints, and statistical testing or decision duality supplies a separate optimality check. A complete output is therefore a procedure together with its assumptions, a validity certificate, an extremizing witness or lower bound, and an explicit account of any remaining gap.

This construction unifies and explains several estimator families without asserting that their sampling experiments are equivalent. It also generates further targets: arbitrary finite multisurvey patterns give kernels, and their approximation, moment, or error-transform problems can be studied by the same machinery. The general sampling and duality ingredients are classical. The value of the present synthesis lies in explicit interfaces, specializations, and the ability to check where an inference ceases to be justified.

The next central target is a certified transfer-and-generation theorem: distinguish positive simulation kernels from signed expectation transfers and approximate experiment reductions, then determine which moment certificates, risks, and calibration guarantees each preserves. The elementary transfer in \cref{prop:transfer} is an input, not a proof of such a general classification. In particular, signed algebraic inversion cannot automatically transfer nonnegative e-values or stopping-time validity.

Several tasks remain. We have not classified all feature families admitting a low-dimensional recovery geometry, solved arbitrary integer-constrained moment problems, constructed a general minimax estimator over all actual populations, or established uniform prediction-error bands over every future horizon. The bounded-batch CLT program is separate and is not imported here as a proved uncertainty-estimation theorem. New-rate-optimality or priority claims require comparison with the cited literature and independent review.

The retained results are already sufficient to support a reusable workflow, but not a universal solution. In particular, a numerical grid optimum is not a continuous dual certificate, a confidence-region projection is not automatically a shortest interval, and exact unbiasedness does not guarantee useful precision.

\appendix
\section{Exact certificate data for the finite tank game}\label{app:certificate}
\Cref{tab:rule} gives a primal rule for \cref{prop:finitegame}. Every displayed probability is exact. Actions not listed at a given observation have probability zero. The current revision checker verifies the clipped rule\textquotesingle s row sums, coverage, every risk constraint, every analytic dual inequality, and equality of primal and dual values using rational arithmetic. The retained legacy certificate verifies the earlier untrimmed rule separately; its non-worst-case risks are not those printed here.
\begin{table}[ht]
\centering\small
\begin{tabular}{cl}
\toprule
Observed $M$ & Reported interval and its probability\\
\midrule
3 & $[3,5]:13/14;\quad [3,7]:3/140;\quad [5,7]:1/20$\\
4 & $[4,6]:23/84;\quad [4,8]:7/15;\quad [4,9]:22/105;\quad [5,10]:1/20$\\
5 & $[5,7]:137/840;\quad [5,12]:211/280;\quad [6,12]:1/12$\\
6 & $[6,9]:41/350;\quad [6,10]:21/100;\quad [6,11]:11/40;$\\
  & $[6,12]:547/1400;\quad [7,12]:1/140$\\
7 & $[7,12]:1$\\
8 & $[8,12]:1$\\
9 & $[9,12]:1$\\
10 & $[10,12]:1$\\
11 & $[11,12]:1$\\
12 & $[12,12]:1$\\
\bottomrule
\end{tabular}
\caption{A support-clipped optimal randomized interval rule in the finite declared action class. It is not asserted to be unique, monotone, or a recommended unrestricted tank procedure.}\label{tab:rule}
\end{table}

\begin{table}[ht]
\centering\small
\begin{tabular}{crr}
\toprule
$N$ & Coverage & Expected interval length\\
\midrule
3 & $19/20$ & $143/70$\\
4 & $19/20$ & $369/112$\\
5 & $19/20$ & $697/140$\\
6 & $19/20$ & $397/80$\\
7 & $19/20$ & $697/140$\\
8 & $19/20$ & $1033/224$\\
9 & $19/20$ & $1369/336$\\
10 & $19/20$ & $1657/480$\\
11 & $19/20$ & $167/60$\\
12 & $19/20$ & $167/80$\\
\bottomrule
\end{tabular}
\caption{Exact coverage and risk for the rule in \cref{tab:rule}.}\label{tab:risks}
\end{table}
\clearpage

\section{Audit ledger and computational scope}\label{app:audit}
This ledger distinguishes verified derivations from assumptions, restrictions, and work not claimed. The audit checks mathematical validity and source positioning; it is not independent peer review.

\begin{longtable}{@{}>{\raggedright\arraybackslash}p{0.25\textwidth}>{\raggedright\arraybackslash}p{0.69\textwidth}@{}}
\toprule
Item & Consolidated status\\
\midrule
\endfirsthead
\toprule Item & Consolidated status\\\midrule\endhead
\bottomrule\endfoot
Sampling units & Retain occurrence indices and unit boundaries. All subset operations are on units, not arbitrary aggregated feature counts.\\[4pt]
Exchangeable differences & Proved on the available horizon. Finite exchangeability is not silently replaced by an infinite mixture model.\\[4pt]
Fixed-sample Good--Turing & Exact expectation is missing mass after $n-1$, not after $n$ observations. Poissonization removes this shift.\\[4pt]
Good--Toulmin inversion & Unbiased for expected future discovery in the Poissonized model; fixed-size extrapolation is a different problem.\\[4pt]
Exponential-integrability gap & Full long-horizon Good--Toulmin has no nonzero exponential moments, even for one species. Direct future-count calibration remains valid. Variance, fourth-moment, and explicitly budgeted truncation alternatives are supplied.\\[4pt]
Dependence & First moments use marginal incidence rates. Covariances require feature intersections. The independent-species cumulant formula is not a batch formula.\\[4pt]
Moment duality & Value duality holds for compact continuous models. Finite dual-certificate attainment at arbitrary boundary points is not claimed.\\[4pt]
Deterministic minimax & Equality is for bounded deterministic error and relaxed linear targets. It is not universal stochastic minimaxity.\\[4pt]
Two-measure theorem & Explicit elementary supporting interpolation proves the envelope. Requires normalization, two specified moments, $a>0$, $t>0$, and relaxed probability measures.\\[4pt]
Even-block extension & Existence uses Karlin--Studden/Pinelis; the endpoint orientation now has an explicit real-rooted Hermite-remainder proof. Odd blocks can switch the order.\\[4pt]
Hill orders & Extended to $(-\infty,2]\cup[3,\infty)$ on $a>0$. Crossing cases in $(2,3)$ have two exhaustive one-parameter support families; actual-population integrality remains separate.\\[4pt]
Population realizability & Atom multiplicities must be integral. Relaxed optimizers can be strictly unattainable. No universal two-level actual-population theorem is claimed.\\[4pt]
Joint calibration & The fixed convex candidate class and finite integrable kernels give a single-level region for $(\mu,Z)$. Arbitrary data-dependent candidate restrictions do not come free. Population-only regions are a specialization.\\[4pt]
Two-node confidence reduction & Valid for regions described only by the two moment coordinates. Additional cumulant kernels require renewed geometry.\\[4pt]
Missing mass & Direct joint calibration is primary; an asymmetric Bernstein alternative is retained. The Esty fixed-sample correction is not omitted or conflated with exact Poisson variance.\\[4pt]
Prediction & Direct target and finite-error joint kernels are separate. Coverage is unconditional for predeclared targets. It is not conditional or automatically uniform over data-selected future windows.\\[4pt]
Chernoff optimization & Sion applies after the perspective change $v=1/\lambda$ on compact convex relaxed classes. Nonconvex integer classes only inherit the conservative minimax inequality.\\[4pt]
Stopping & Fixed-time e-values are not automatically confidence sequences. Sequential validity requires the stated supermartingale or error-spending construction.\\[4pt]
Affine risk program & Exact algebraic risk plus a relaxed uniform bound. Does not prove minimaxity against all estimators on actual populations.\\[4pt]
Tank upper endpoint & The strict $>$ acceptance condition gives pointwise minimality; monotonicity is proved by nested acceptance sets with $F_N(m)=1$ for $m\ge N$.\\[4pt]
Finite decision game & $697/140$ is proved for the explicit finite parameter/action class with randomization and length $b-a$. The primal rule is clipped to $[M,12]$, reducing four non-worst-case risks. The unchanged analytic dual proves the same optimum.\\[4pt]
Code verification & Original checks rerun; a separate rational-certificate checker and symbolic/numerical additions included. No formal proof assistant verification or broad software validation is asserted.\\[4pt]
Novelty & Classical components are attributed. The draft makes no blanket first-unification, universal-compiler, or new-minimax-rate claim.\\
\end{longtable}

\subsection{Reproducibility}
The source bundle includes the original three verification modules, their rerun outputs, the historical finite-game certificate, an independent audit script, and a separate revision checker for the clipped certificate and the added results. Historical and current certificates are retained under distinct names; the manuscript tables use the clipped certificate. The checks have distinct meanings:
\begin{center}\small
\begin{tabular}{p{0.47\textwidth}rp{0.32\textwidth}}
\toprule
Check family & Count & Arithmetic / interpretation\\
\midrule
Sampling identities & 402 & Exact rational finite examples\\
Recovery-geometry checks & 7081 & Mostly floating-point; four rational checks\\
Occupancy cumulant directions & 576 & Enumerated finite laws, floating-point\\
Tank coverage cases & 465 & Exact integer/rational comparisons\\
Tank minimality witnesses & 225 & Exact integer/rational comparisons\\
Prediction-error transforms & 18 & Enumerated finite laws, floating-point\\
Independent finite-game constraints & 593 & Exact rational primal--dual verification\\
Operator identities & 3 & Symbolic simplification\\
Additional $q\ge3$ envelopes & 250 & Floating-point checks\\
\bottomrule
\end{tabular}
\end{center}
These retained-suite counts are not a substitute for proofs, and some belong to overlapping test runs. The current revision additionally checks 603 rational constraints for the clipped game, 360 joint-transform comparisons, three symbolic zero-rate limits, 600 Hill-order comparisons including negative orders, and 36 truncation probabilities. The recovery run's largest moment error was approximately $1.65\cdot10^{-14}$ and its largest grid-LP discrepancy approximately $2.38\cdot10^{-11}$. The cumulant identity errors were below $8\cdot10^{-15}$ in the retained cases. These are empirical numerical residuals, not rigorous floating-point error bounds.

The revision checker additionally verifies the direct joint transforms, their zero-rate limits, fourth moments and the non-improvement example, negative-order envelopes, the two intermediate-Hill support families against finite-grid linear programs, the exact truncation transform, and the perspective identity. These extra numerical checks are corroboration, not global optimization certificates; their counts and residuals are recorded in \texttt{checks/revision\_v02.json}.

The original exponential calibrations were conservative in the 24 enumerated cases: nominal $95\%$ regions had minimum observed exact-law coverage above $99.7\%$ for the tested mixture. This is evidence about those examples only, and it is not an efficiency theorem.

The two-node solver uses double precision and tolerance-based boundary recognition. It is intended for moderately scaled inputs, not near-singular moment data or arbitrarily large $tb$. Production certification would require interval arithmetic, carefully bracketed roots, and continuum inequality certificates. No external dataset or simulation is used to justify a theorem.

\begingroup\small
\begin{thebibliography}{99}
\setlength{\itemsep}{3pt}\setlength{\parskip}{0pt}
\bibitem{BenHamou2017}
A. Ben-Hamou, S. Boucheron, and M. I. Ohannessian.
Concentration inequalities in the infinite urn scheme for occupancy counts and the missing mass, with applications.
\emph{Bernoulli} \textbf{23}(1) (2017), 249--287.
\href{https://doi.org/10.3150/15-BEJ743}{doi:10.3150/15-BEJ743}.

\bibitem{BertsimasPopescu2005}
D. Bertsimas and I. Popescu.
Optimal inequalities in probability theory: A convex optimization approach.
\emph{SIAM Journal on Optimization} \textbf{15}(3) (2005), 780--804.
\href{https://doi.org/10.1137/S1052623401399903}{doi:10.1137/S1052623401399903}.

\bibitem{Boucheron2013}
S. Boucheron, G. Lugosi, and P. Massart.
\emph{Concentration Inequalities: A Nonasymptotic Theory of Independence}.
Oxford University Press, 2013.
\href{https://doi.org/10.1093/acprof:oso/9780199535255.001.0001}{doi:10.1093/acprof:oso/9780199535255.001.0001}.

\bibitem{Chao1984}
A. Chao.
Nonparametric estimation of the number of classes in a population.
\emph{Scandinavian Journal of Statistics} \textbf{11}(4) (1984), 265--270.

\bibitem{ChaoJost2015}
A. Chao and L. Jost.
Estimating diversity and entropy profiles via discovery rates of new species.
\emph{Methods in Ecology and Evolution} \textbf{6}(8) (2015), 873--882.
\href{https://doi.org/10.1111/2041-210X.12349}{doi:10.1111/2041-210X.12349}.

\bibitem{DaleySmith2016}
T. Daley and A. D. Smith.
Better lower bounds for missing species: improved non-parametric moment-based estimation for large experiments.
Preprint (2016; revised March 4, 2019), \href{https://arxiv.org/abs/1605.03294v3}{arXiv:1605.03294v3}.

\bibitem{Donoho1994}
D. L. Donoho.
Statistical estimation and optimal recovery.
\emph{Annals of Statistics} \textbf{22}(1) (1994), 238--270.
\href{https://doi.org/10.1214/aos/1176325367}{doi:10.1214/aos/1176325367}.

\bibitem{EfronThisted1976}
B. Efron and R. Thisted.
Estimating the number of unseen species: How many words did Shakespeare know?
\emph{Biometrika} \textbf{63}(3) (1976), 435--447.
\href{https://doi.org/10.1093/biomet/63.3.435}{doi:10.1093/biomet/63.3.435}.

\bibitem{Eriksson2026}
E. Eriksson.
The unseen species problem revisited.
Preprint (2026), \href{https://arxiv.org/abs/2602.08769v4}{arXiv:2602.08769v4}, September 9, 2026.

\bibitem{Esty1983}
W. W. Esty.
A normal limit law for a nonparametric estimator of the coverage of a random sample.
\emph{Annals of Statistics} \textbf{11}(3) (1983), 905--912.
\href{https://doi.org/10.1214/aos/1176346256}{doi:10.1214/aos/1176346256}.

\bibitem{Good1953}
I. J. Good.
The population frequencies of species and the estimation of population parameters.
\emph{Biometrika} \textbf{40}(3--4) (1953), 237--264.
\href{https://doi.org/10.1093/biomet/40.3-4.237}{doi:10.1093/biomet/40.3-4.237}.

\bibitem{GoodToulmin1956}
I. J. Good and G. H. Toulmin.
The number of new species, and the increase in population coverage, when a sample is increased.
\emph{Biometrika} \textbf{43}(1--2) (1956), 45--63.
\href{https://doi.org/10.1093/biomet/43.1-2.45}{doi:10.1093/biomet/43.1-2.45}.

\bibitem{Harris1959}
B. Harris.
Determining bounds on integrals with applications to cataloging problems.
\emph{Annals of Mathematical Statistics} \textbf{30}(2) (1959), 521--548.
\href{https://doi.org/10.1214/aoms/1177706266}{doi:10.1214/aoms/1177706266}.

\bibitem{Howard2021}
S. R. Howard, A. Ramdas, J. McAuliffe, and J. Sekhon.
Time-uniform, nonparametric, nonasymptotic confidence sequences.
\emph{Annals of Statistics} \textbf{49}(2) (2021), 1055--1080.
\href{https://doi.org/10.1214/20-AOS1991}{doi:10.1214/20-AOS1991}.

\bibitem{KarlinStudden1966}
S. Karlin and W. J. Studden.
\emph{Tchebycheff Systems: With Applications in Analysis and Statistics}.
Pure and Applied Mathematics, vol.~15. Interscience Publishers, New York, 1966.

\bibitem{LastPenrose2017}
G. Last and M. Penrose.
\emph{Lectures on the Poisson Process}.
Institute of Mathematical Statistics Textbooks, vol.~7. Cambridge University Press, 2017.
\href{https://doi.org/10.1017/9781316104477}{doi:10.1017/9781316104477}.

\bibitem{ONeill2025}
B. O'Neill.
The distribution of order statistics under sampling without replacement.
\emph{Journal of Statistical Theory and Applications} \textbf{24}(3) (2025), 663--698.
\href{https://doi.org/10.1007/s44199-025-00125-y}{doi:10.1007/s44199-025-00125-y}.

\bibitem{Orlitsky2016}
A. Orlitsky, A. T. Suresh, and Y. Wu.
Optimal prediction of the number of unseen species.
\emph{Proceedings of the National Academy of Sciences} \textbf{113}(47) (2016), 13283--13288.
\href{https://doi.org/10.1073/pnas.1607774113}{doi:10.1073/pnas.1607774113}.

\bibitem{Pinelis2011}
I. Pinelis.
Tchebycheff systems and extremal problems for generalized moments: a brief survey.
Preprint (2011), \href{https://arxiv.org/abs/1107.3493}{arXiv:1107.3493}.

\bibitem{PolyanskiyWu2019}
Y. Polyanskiy and Y. Wu.
Dualizing Le Cam's method for functional estimation, with applications to estimating the unseens.
Preprint (2019; updated versions), \href{https://arxiv.org/abs/1902.05616}{arXiv:1902.05616}.

\bibitem{PolyanskiyWu2026}
Y. Polyanskiy and Y. Wu.
Dualizing Le Cam's method for functional estimation I: General theory.
\emph{Annals of Statistics} \textbf{54}(1) (2026), 1--24.
\href{https://doi.org/10.1214/25-AOS2498}{doi:10.1214/25-AOS2498}.

\bibitem{Sion1958}
M. Sion.
On general minimax theorems.
\emph{Pacific Journal of Mathematics} \textbf{8}(1) (1958), 171--176.
\href{https://doi.org/10.2140/pjm.1958.8.171}{doi:10.2140/pjm.1958.8.171}.

\bibitem{Wasserman2020}
L. Wasserman, A. Ramdas, and S. Balakrishnan.
Universal inference.
\emph{Proceedings of the National Academy of Sciences} \textbf{117}(29) (2020), 16880--16890.
\href{https://doi.org/10.1073/pnas.1922664117}{doi:10.1073/pnas.1922664117}.
\bibitem{WuYang2019}
Y. Wu and P. Yang.
Chebyshev polynomials, moment matching, and optimal estimation of the unseen.
\emph{Annals of Statistics} \textbf{47}(2) (2019), 857--883.
\href{https://doi.org/10.1214/17-AOS1665}{doi:10.1214/17-AOS1665}.

\bibitem{ZhangZhang2009}
C.-H. Zhang and Z. Zhang.
Asymptotic normality of a nonparametric estimator of sample coverage.
\emph{Annals of Statistics} \textbf{37}(5A) (2009), 2582--2595.
\href{https://doi.org/10.1214/08-AOS658}{doi:10.1214/08-AOS658}.
\end{thebibliography}
\endgroup
\end{document}
